\documentclass[11pt]{article}
\usepackage{fontspec}
\usepackage{xeCJK}
\setmainfont{Noto Serif}
\setCJKmainfont{Noto Serif CJK JP}
\setCJKsansfont{Noto Sans CJK JP}
\XeTeXlinebreaklocale "ja"
\XeTeXlinebreakskip=0pt plus 0.1pt

\usepackage{amsmath,amssymb,mathtools,array,booktabs,pdflscape,adjustbox,diagbox}
\usepackage{geometry}
\geometry{a4paper,margin=1.45cm}
\setlength{\parindent}{1.25em}
\setlength{\parskip}{0pt}
\makeatletter
\renewcommand\footnotesize{\@setfontsize\footnotesize{7.5}{8.7}\selectfont}
\makeatother
\allowdisplaybreaks
\setlength{\emergencystretch}{12em}
\sloppy
\hbadness=10000
\vbadness=10000
\hfuzz=2pt
\vfuzz=10pt
\raggedbottom
\providecommand{\widebar}[1]{\overline{#1}}
\providecommand{\mR}{\mathfrak{R}}
\providecommand{\mK}{\mathfrak{K}}
\providecommand{\mL}{\mathfrak{L}}
\providecommand{\mS}{\mathfrak{S}}
\providecommand{\mT}{\mathfrak{T}}
\providecommand{\mM}{\mathfrak{M}}
\providecommand{\mN}{\mathfrak{N}}
\providecommand{\mA}{\mathfrak{A}}
\providecommand{\mB}{\mathfrak{B}}
\providecommand{\mC}{\mathfrak{C}}
\providecommand{\mD}{\mathfrak{D}}
\providecommand{\mE}{\mathfrak{E}}
\providecommand{\mF}{\mathfrak{F}}
\providecommand{\mG}{\mathfrak{G}}
\providecommand{\mH}{\mathfrak{H}}
\providecommand{\mU}{\mathfrak{U}}
\providecommand{\mV}{\mathfrak{V}}
\providecommand{\mbar}[1]{\overline{#1}}
\providecommand{\ma}{\mathfrak{a}}
\providecommand{\mb}{\mathfrak{b}}
\providecommand{\mc}{\mathfrak{c}}
\providecommand{\md}{\mathfrak{d}}
\providecommand{\me}{\mathfrak{e}}
\providecommand{\mf}{\mathfrak{f}}
\providecommand{\mm}{\mathfrak{m}}
\providecommand{\mn}{\mathfrak{n}}
\providecommand{\mo}{\mathfrak{o}}
\providecommand{\mpideal}{\mathfrak{p}}
\providecommand{\mq}{\mathfrak{q}}
\providecommand{\mt}{\mathfrak{t}}
\providecommand{\mv}{\mathfrak{v}}
\providecommand{\frakp}{\mathfrak p}
\providecommand{\frakq}{\mathfrak q}
\providecommand{\modu}[1]{\;\operatorname{mod} #1}

\begin{document}

\section*{30. 代数的数体および関数体におけるイデアル論の抽象的構成}
\noindent\emph{Math. Ann. 96 (1927), pp. 26--61.}

以下では、代数的数体のすべての整数のイデアル論と一致するイデアル論をもつすべての環、すなわちそのイデアルが素イデアルの冪の積として一意に表される環の抽象的特徴づけを与える。このような環のよく知られた例には、一変数の代数的関数体のすべての整数からなる環も含まれ、さらに一般に、Kronecker に従って最高次元のイデアルに制限する、すなわちある商環である関数領域へ移るなら、複数の不定元の場合も含まれる。

基礎に置く可換環について、このイデアル論と完全に同値な公理は次のものである。\footnote{イデアル論の基本概念の定義については、例えば E. Noether, \emph{Idealtheorie in Ringbereichen}, Math. Ann. 83 (1921), pp. 24--66（以下では \emph{イデアル論}として引用）を参照。なお、本論文、とくに §§1, 4, 5 でも、最も本質的な概念を簡潔に定式化している。完全を期して、本論文の直接の目的には必要でない材料もいくらか含めた。}

\begin{description}
\item[I. 約イデアル鎖条件。] 各イデアルが直前のものの真の約イデアルであるような任意のイデアル鎖は、有限段階の後に停止する。言い換えれば、イデアルの任意の約鎖には、以後すべてのイデアルが等しくなる添字がある。
\item[II. 任意の非零イデアルを法とする倍イデアル鎖条件。] 零イデアルと異なる固定イデアルの約イデアルだけからなり、各イデアルが直前のものの真の倍イデアルであるような任意のイデアル鎖は、有限段階の後に停止する。
\item[III.] 乗法の単位元の存在。
\item[IV.] 零因子をもたない環。
\item[V. 商体における整閉性。] 商体の元で環に関して整なものは、すべてその環に属する。
\end{description}

これらの公理から、望まれるものに至るまで、次第に制限されたイデアル論を段階的に展開する。逆に、これらすべての公理が実際にそこで満たされることを示す。

約イデアル鎖条件 I からは、私が \emph{イデアル論} §4 で示したように、イデアルが、相異なる素イデアルに属する有限個の準素イデアルの最小公倍数として表されることが従う。この証明をここで簡単に想起する（§6）。その理由は、イデアル商の概念を用いることで簡単化できるからであり、また一つの見落としを訂正したいからでもある。元の証明は一箇所で、仮定されていない乗法単位元の存在を用いていた。\footnote{P. Urysohn がこの見落としを私に指摘してくれた。私自身による証明の変更は、ここに与える E. Artin によるものよりやや煩雑であった。彼はすでに以前この欠落を自分で埋めており、また私とは独立に、イデアル商による簡単化にも気づいていた。「簡約表示」の導入を避けるさらなる簡単化は、W. Krull, \emph{Algebraische Theorie der Ringe III}, Math. Ann. 92 (1924), pp. 181--213, 定理 I の関連する問題から取った。}

倍イデアル鎖条件の仮定は（§7）、零イデアルと異なる任意のイデアルを法とする剰余類環において、素イデアルが、全元を含む単位イデアルを除いて、真の約イデアルをもたないという効果をもつ。従って、これらのイデアルの準素成分は、単位イデアルに属するものを除いて一意に定まる。やや弱い形ではこの結果はすでに Masazo Sono に見られる。\footnote{Masazo Sono, \emph{On the Reduction of Ideals}, Memoirs of the College of Science, Kyoto University, Series A, 7 (1924), pp. 191--204.} 彼の、剰余類環に組成列が存在するという仮定は、この環における「二重鎖条件」と同値であることを末尾（§10）で簡単に示す。二重鎖条件とは、零イデアルと異なるイデアルへの制限なしに、約イデアル鎖条件と倍イデアル鎖条件の双方が成り立つことをいう。

さらに公理 III、すなわち単位元の存在が満たされるなら、単位イデアルは随伴素イデアルとして現れない。従って準素成分は一意に定まり、互いに二つずつ互いに素であり、その最小公倍数は積に等しい。公理 I, II, III は代数的数体のすべての有限オーダーに対して成り立つ。すでに Dedekind（Dirichlet--Dedekind, \emph{Zahlentheorie}, 第3版, 第11補遺, §172）は、証明を完全には詳述しないまま、イデアルを互いに素な一種イデアル（準素成分）の積として一意に表すことを述べていた。

仮定 IV、すなわち零因子がないことからは、零イデアルでも単位イデアルでもないイデアルの異なる冪はすべて相異なること、および商体が存在することが従う。最後に、商体における整閉性の公理 V も満たされるなら、IV により一意に定まる準素イデアルは素イデアルの冪となり、通常の因数分解定理が得られる（§8）。この最後の証明は、Dedekind が数体の場合に（第3版, §172）すでに整閉性から引き出した帰結に基づく。すなわち、真に分数的な元は、分子の二乗さえ分母で割り切れないように、整数の商として表せるという帰結である。後の叙述では、この帰結は同じく整閉性の帰結として、はるかに見通しの悪い一般化 Gauss 定理、あるいはそれよりやや強い対応する加群定理によって置き換えられる。これらはいずれも基底元を前提として適用される定理であるが、ここではイデアルそのものを直接扱う。

逆方向（§9）では、公理 V 以外の公理はほとんど直接に従う。V は、真に分数的な元が、分子のいかなる冪も分母で割り切れないように常に表されるという証明から従い、これはもともと V から導かれた帰結を強めるものである。この証明が成功するのは、イデアル表示から通常の帰結を引き出した後に限られる。すなわち、固定イデアルを法とする主イデアルによる表示と、分数イデアルの理論、つまり体におけるイデアル商の理論である。素イデアルの冪による表示から整閉性が従うことは Dedekind にもすでに知られていた。第3版 §172, p. 522 下部の注記がそれを示している。

公理 I から V は、一般には異なる四つの分解、すなわち二つずつ互いに素な最小イデアルへの分解、互いに素なイデアルへの分解、最大準素成分への分解、既約イデアルへの分解が一致することを含意する。逆に、この一致と公理 I--IV から整閉性が従う。零因子をもつ環では、残りの公理を満たしても、V を商環における整閉性として定義した場合でさえ、必ずしもこれらの分解が一致するとは限らない。有限オーダーのあるイデアルを法とする剰余類環の例がこれを示す（§9）。

§1 では、環に関して整な量の理論を、整閉性から Dedekind が得た帰結まで概観する。§2 では、鎖条件が、その環を係数および乗数領域とする線型形式からなる加群へどのように移されるかを示す。\footnote{この形の加群定理は Prüfer（\emph{Neue Begründung der algebraischen Zahlentheorie}, §2, Math. Ann. 94 (1924), pp. 198--243）によるもので、基底定理を移す通常の方法より見通しがよい。} そこから（§3）、第一種の代数的拡大体の整元へ移ると、公理 I から IV は任意のオーダーについて保たれ、すべての整元からなるオーダーでは V も成り立つ、という結論を得る。この事実は、もちろん代数的数体では周知であるが、冒頭で述べた関数体を従属させることを可能にする。特に、複数の不定元の整多項式から導かれる関数領域が公理を満たすことの簡単な証明により、Kronecker のイデアル論は完全に包含される。\footnote{違いは判別式定理にだけ現れる。素数を法とする関数領域の剰余類環が不完全体になり、従って第二種拡大が起こり得るからである。}

後で展開されるイデアル論は、分類だけに用いられる §§2, 3 を前提としない。§4 と §5 では、公理 I から V に依存しない基礎を与える。これにより、理論のどの部分が有限性仮定に依存しないかが、元の基礎づけよりも鋭く浮き出る。

\subsection*{§1. 整量の理論}

可換環\footnote{よく知られているように、環は、集合上に等号関係が与えられ、さらに二つの演算、加法と乗法が与えられて通常の法則を満たすとき定義される。加法は結合的・可換的で一意に可逆であり、乗法は結合的で、環が可換であるときは可換的であり、加法に関して分配的である。基礎となる等号の定義が集合論的同一性でない場合には、すべての部分系で同じ等号の定義が保たれるように、その部分系（部分環、加群、イデアル）は任意の元とともにそれに等しいすべての元を含まなければならない。同様に、すべての拡大では新しい等号の定義が古いものを含むと仮定しなければならない。「有限個の元」「一意の対応」などのすべての概念は、この等号定義の意味で理解される。例えば「一つだけの元がある」とは、等号を除いて一つだけであることを意味する。なお、等しい元を一つの類にまとめ、これらの類を新しい元として導入することで、等号から集合論的同一性へ常に移れる。ここに述べた等号定義に関する注意は R. Hölzer による。} $T$ が与えられているとする。これは零因子をもたず、乗法の単位元 $e$ をもつ（公理 III, IV）。$T$ の中には、単位元を含む固定された部分環 $R$ が区別されている。加群、オーダー、整という概念は、常にこの固定環 $R$ に関するものとし、後には簡単のため記号から省く。\footnote{整量の理論は、零因子を許しても、それに代えて $R$ における約イデアル鎖条件を仮定すれば有効であり、§2 によりこの仮定は補題の証明も可能にする。この形の整量の理論は後で使わないので、注で指摘するにとどめる。ここで問題となる定義はすべて、零因子が存在しても有効である。その多くは、よく知られ容易に見られるように、さらに弱い仮定しか必要としない。保たれないのは Dedekind の「数は包絡をもつとき整である」という定式化である。零因子があると、二つの有限加群の商が有限であるとは限らないからである。したがってここではオーダーを加群商としてでなく直接定義する。} $R$ の元はラテン文字で、$T$ の一般の元はギリシア文字で表す。

1. $T$ の元からなる系 $M$ は、通常どおり、二つの元 $\alpha,\beta$ とともにその差を含み、$\alpha$ とともに $r\alpha$ も含むとき、$R$-加群、略して加群と呼ばれる。ただし $r$ は $R$ の任意の元である。$M$ が $N$ で割り切れるとは
\[
        M\equiv 0 \pmod N
\]
で表し、$M$ を倍加群、$N$ を約加群と呼ぶ。$R$-加群で、そのすべての元が $R$ に属するものを $R$ のイデアルと呼ぶ。

明らかに、$T$ の任意個の加群の共通部分、すなわち最小公倍数 $[\ldots,M_i,\ldots]$ は再び加群である。$T$ 自身も加群である。したがって $T$ の元からなる任意の系 $\Sigma$ から導かれる加群は、$\Sigma$ を含むすべての加群の共通部分として一意に存在する。任意の加群系の和、すなわち最大公約数 $(\ldots,M_i,\ldots)$ は、それらの和集合から導かれる加群である。二つの加群の積 $MN$ は、$\mu$ が $M$ のすべての元を、$\nu$ が $N$ のすべての元を走るときの元 $\mu\nu$ から導かれる加群として定義される。従って積は、環の元について成り立つため、結合律と可換律を満たす。さらに $T$ における分配律から、加群についても分配律
\[
        (\ldots,M_i,\ldots)N=(\ldots,M_iN,\ldots)
\]
が従う。まさに $T$ におけるこの法則によって、両辺はいずれもすべての元 $\mu_i\nu$ から導かれる加群だからである。

$R$ 自身が加群であるから、$R$ が乗数領域であるという条件は $RM\equiv0\pmod M$ と表せる。仮定により $R$ は単位元を含むので、$M\equiv0\pmod{RM}$ も成り立ち、従って $M=RM$ である。

加群 $M$ は、それを導く有限個の元からなる系 $\Sigma$ が少なくとも一つ存在するとき有限と呼ばれる。$\Sigma$ の元 $\mu_1,\ldots,\mu_n$ は
\[
        M=(\mu_1,\ldots,\mu_n)
\]
の加群基底と呼ばれる。二つ、従って有限個の有限加群の和および積は再び有限である。それぞれ基底元の和集合、あるいは基底元の積から導かれるからである。後者の主張は、$M$ のすべての元 $\mu$ が基底によって $r_1\mu_1+\cdots+r_n\mu_n$ と表されること、および $T$ における乗法の可換律から従う。

2. $T$ の中にある $R$ の拡大環、すなわち $R$ のすべての元を含む環を $R$-オーダー、略してオーダーと呼ぶ。$T$ における任意個のオーダーの共通部分は再びオーダーであり、$T$ もまたオーダーである。したがって $T$ の元からなる任意の系 $\Sigma$ から導かれるオーダーは一意に存在する。オーダーの和および積は加群の場合と同様に定義され、特に任意のオーダーについて $O^2=O$ である。

すべてのオーダーは同時に $R$-加群である。オーダーは有限加群であるとき有限と呼ばれる。和と積は加群の和と積を作ることで得られるので、1 により有限オーダーの和と積は再び有限オーダーである。さらに次の推移律が成り立つ。$A$ が有限 $R$-オーダーで、$B$ が有限 $A$-オーダーなら、$B$ は有限 $R$-オーダーでもある。実際、$B$ は同時に $R$-オーダーであり、従って $R$-加群である。$\alpha_1,\ldots,\alpha_m$ が $R$ に関する $A$ の加群基底で、$\beta_1,\ldots,\beta_n$ が $A$ に関する $B$ の加群基底なら、積 $\alpha_i\beta_j$ は明らかに $R$ に関する $B$ の加群基底をなす。

3. $T$ の元 $\alpha$ は、$\alpha$ から導かれるオーダー $R_\alpha$ が有限であるとき、$R$ に関して整、略して整と呼ばれる。同値に、加群鎖
\[
        U_n=(\alpha^0,\alpha,\ldots,\alpha^{n-1})
\]
が有限段階の後に停止し、$U_n=U_{n+1}=\cdots$ となるときである。

この定義には、4. で用いる第二の定式化もある。$T$ の元 $\alpha$ は、零加群と異なる主加群 $C$、すなわちある $\gamma\ne0$ から導かれる $C$ が少なくとも一つ存在し、オーダー $R_\alpha$ と $C$ の積が有限加群になるとき、整と呼ばれる。言い換えれば、加群鎖 $CU_n$ が有限段階の後に停止することである。\footnote{$R$ に零因子を許す場合には、正則な主加群の存在を要求しなければならない。すなわち $\gamma$ は零因子でない正則元であると仮定しなければならず、零因子のない環ではこれは単に $\gamma\ne0$ である。}

最後に、この定義は通常の定式化とも同一である。$T$ の元 $\alpha$ は、$R$ の元 $r_i$ を用いて少なくとも一つの方程式
\[
        \alpha^n+r_1\alpha^{n-1}+\cdots+r_n=0
\]
を満たすとき整と呼ばれる。

これらの定義は実際に同一である。$R_\alpha$ は $\alpha$ のすべての冪から導かれる加群と一致するから、$R_\alpha$ が有限であるとき、その冪の中から常に加群基底を選ぶことができる。そのような基底が、例えば $e=\alpha^0,\alpha,\ldots,\alpha^{n-1}$ で与えられるなら、これは $U_n$ の鎖が $U_n$ で停止することを意味する。逆に $U_n=U_{n+1}=\cdots$ はこの基底を与える。さらに $U_n$ の鎖の停止は、方程式
\[
        \alpha^n+r_1\alpha^{n-1}+\cdots+r_n=0
\]
の存在とも同一である。$R_\alpha$ が有限であれば $CR_\alpha$ も有限であり、従って鎖 $CU_n$ は停止する。逆に $CR_\alpha$ の有限性、したがって $CU_n$ の停止から、$U_n$ も停止し、従って $R_\alpha$ が有限であることが従う。実際、$C$ は零加群と異なる主加群であると仮定されているので、$CU_n=CU_{n+1}$ は $U_n=U_{n+1}$ も与える。

整量の環性と整依存の推移律は、次の補題に基づく。

\medskip
\noindent\textbf{補題。} $O$ を $T$ の有限オーダーとし、$\alpha$ を $O$ の任意の元とする。このとき $\alpha$ から導かれるオーダー $R_\alpha$ は有限である。言い換えれば、有限オーダーのすべての元は整である。

証明は通常の議論による。$\omega_1,\ldots,\omega_n$ を $O$ の加群基底とする。$O$ は環であるから、$\alpha\omega_i$ は $O$ に属する。従って
\[
        \alpha\omega_i=r_{i1}\omega_1+\cdots+r_{in}\omega_n,
\]
したがって
\[
        \det(r_{ij}-e_{ij}\alpha)=0
\]
である。ただし $i\ne j$ なら $e_{ij}=0$、$e_{ii}=e$ である。これにより $\alpha$ が整であることが示される。行列式の消滅には、$T$、従って $O$ が零因子をもたないという仮定を用いる。\footnote{$R$ に約イデアル鎖条件を仮定しつつ零因子を許す場合、この補題は §2 の加群定理から直ちに従う。その定理によれば鎖条件は $O$ からなる加群へ移る。この証明も、数体については Dedekind（第4版, §173, II）に由来し、文献における鎖条件の最初の応用である。4. の下で与える帰結では、Dedekind は代数的数体におけるイデアルを法とする剰余類の数が有限であるという、より特殊な事実を用いている。}

$T$ におけるすべての整量からなる系 $G$ はオーダーをなす。$R$ の元は、単位元を基底とする有限 $R$-加群として $R$ が有限であるため整であり、したがって $G$ は $R$ を含む。また $G$ は環性ももつ。任意の二つの整量 $\alpha,\beta$ から導かれるオーダー $R_{\alpha,\beta}$ は積 $R_\alpha R_\beta$ に等しく、従って有限である。ゆえに補題により、$\alpha+\beta$ と $\alpha\beta$ は有限オーダーの元として整である。

\medskip
\noindent\textbf{整依存の推移律。} $O$ が整量からなるオーダーであるなら、$O$ に関して整な $T$ のすべての元は $R$ に関しても整である。定義により $\alpha$ は方程式
\[
        \alpha^m+\omega_1\alpha^{m-1}+\cdots+\omega_m=0
\]
を満たし、従って $\omega_1,\ldots,\omega_m$ から導かれるオーダー $O'=R_{\omega_1,\ldots,\omega_m}$ に関しても整である。この $O'$ は積 $R_{\omega_1}\cdots R_{\omega_m}$ として有限である。2. で与えた推移律により、$\alpha$ から導かれる有限 $O'$-オーダー $O'_\alpha$ は従って有限 $R$-オーダーでもあり、$\alpha$ は有限オーダーの元として補題により整である。

4. 環 $R$ は、$R$ に関して整な $T$ のすべての元が $R$ に属するとき、$T$ において整閉であると呼ばれる。3. の整量定義の第二の定式化によれば、$T$ の元 $\alpha$ が $R$ に属するための必要十分条件は、零加群と異なる主加群 $C$ が少なくとも一つ存在して $CR_\alpha$ が有限加群となることである。\footnote{$R$ に零因子を許す場合、$C$ は再び正則主加群と仮定しなければならない。同様に、帰結 I の元 $c$ も正則と仮定しなければならない。}

整閉環の概念から Dedekind（第3版, §172）は、これをイデアル論に用いることを可能にする二つの重要な帰結を引き出した。

\medskip
\noindent\textbf{Dedekind の帰結 I。} $R$ が $T$ において整閉であり、さらに $R$ でイデアルの約イデアル鎖条件（公理 I）が成り立つとする。$\beta$ が $T$ の非整元で、$c\ne0$ が $R$ の元なら、元の列
\[
        c,\;c\beta,\ldots,c\beta^\nu,\ldots
\]
において $c,\ldots,c\beta^\rho$ は整であり、残りはすべて非整であるような指数 $\rho>0$ が存在する。

もしすべての元 $c\beta^\nu$ が整であれば、$R$ の整閉性によりそれらは $R$ に属する。すると、$C$ を $c$ から導かれる加群、$B_\nu$ を $1,\beta,\ldots,\beta^\nu$ から導かれる加群とすると、加群鎖 $C,CB_1,\ldots,CB_\nu,\ldots$ はイデアル鎖 $C_\nu=(c,c\beta,\ldots,c\beta^\nu)$ になり、仮定により停止する。しかしそのとき、仮定に反して $R_\beta$ は有限となり、$\beta$ は整となる。従って、元 $c\beta^\nu$ のいくつかは非整である。さらに、一つの非整元があれば、それ以後のすべての元は非整である。なぜなら、$c\beta^r$ が整で、従って $R$ に属するなら、$\rho<r$ に対して元 $c\beta^\rho$ は
\[
        (c\beta^\rho)^r-c^{\,r-\rho}(c\beta^r)^\rho=0
\]
を満たし、従って整だからである。従って $c\beta^{\rho+1}$ が最初の非整元であれば、$c$ は整であるから $\rho>0$ であり、これが求める指数である。

拡大環 $T$ を $R$ の商体、すなわちすべての元の対を添加して得られる体\footnote{Steinitz のよく知られた定式化、J. f. M. 137 において。$R$ に零因子を許す場合には、分母が $R$ の正則元である商対すべてを添加して、商体の代わりに商環を用いなければならない。この場合、Dedekind の帰結 II はもはや有効でない。$n$ は一般に正則元ではないからである。}に特殊化すると、次が得られる。

\medskip
\noindent\textbf{Dedekind の帰結 II。} $R$ がその商体 $K$ の中で整閉であり、$R$ でイデアルの約イデアル鎖条件が成り立つとする。$\beta$ が $K$ の非整元なら、$\beta$ は $R$ の元の商
\[
        \beta=m/n
\]
として表され、しかも $m^2/n$ も非整となるようにできる。

$\beta=b/c$ と置き、$c\ne0$ とする。このとき列 $c,c\beta=b,c\beta^2,\ldots$ の最初の二つの元は確かに整であるから、帰結 I の指数を $\rho$ とすれば $\rho>1$ である。そこで
\[
        m=c\beta^\rho,\qquad n=c\beta^{\rho-1}
\]
と置けば、$m/n=\beta$ であり、$m^2/n=c\beta^{\rho+1}$ は非整である。

整閉性の推移律は次の形をとる。$R$ を $T$ の任意の環とし、$G$ を $R$ に関して整な $T$ のすべての量の系とする。このとき $G$ は、$G$ に関して整な $T$ のすべての量からもなる。実際、$G$ に関して整な任意の量は、3. の推移律により $R$ に関しても整であり、従って $G$ に属する。

\subsection*{§2. 有限加群領域における鎖条件}

鎖条件を移す加群定理は、ここでの応用では拡大環の加群についてだけ現れる。証明は一般の加群領域の場合でも同じであるから、そのような領域を出発点とする。

元 $\alpha,\beta,\ldots$ からなる系 $M$ は、環 $R$ に関して、$M$ の元へ一意に導く二つの演算、すなわち元の加法的結合と $R$ の元による乗法が与えられ、$M$ が加法について Abel 群であり、乗法が結合律を満たし、分配律が二つの形で成り立つとき、加群領域と呼ばれる。\footnote{\emph{イデアル論}, §9 を参照。}

加群領域については、§1, 1. で与えた加群定義のうち、乗法に関係しないものはすべて明らかに有効である。特に、$M$ は、$M$ の有限個の元 $\xi_1,\ldots,\xi_k$ が存在し、$M$ が $\xi_1,\ldots,\xi_k$ から導かれる加群、従って $R$ が単位元をもつ場合にはすべての線型形式 $r_1\xi_1+\cdots+r_k\xi_k$ の系に等しいとき、有限加群領域と呼ばれる。そうでない場合には追加項 $n_i\xi_i$ が現れる。

\medskip
\noindent\textbf{加群定理。} $M$ を、単位元をもつ可換環 $R$ に関する有限加群領域とし、$R$ においてイデアルに対する約イデアル鎖条件、または倍イデアル鎖条件が成り立つとする。このとき $M$ では、加群に対する約加群鎖条件、または倍加群鎖条件が成り立つ。

証明は、$M$ の各加群に対して $R$ の有限個のイデアルからなる系を一意に対応させ、加群の真の約加群または真の倍加群に対して、そのたびに少なくとも一つの真の約イデアルまたは倍イデアルが対応するようにすることに基づく。

$M$ の元が長さ $i$ をもつとは、それが少なくとも一つの仕方で $\xi_1,\ldots,\xi_i$ の線型形式として書けるが、$\xi_1,\ldots,\xi_{i-1}$ の線型形式としては表されないことをいう。$M$ の任意の加群 $W$ に、$R$ の $k$ 個のイデアル $\mathfrak a_1,\ldots,\mathfrak a_k$ を割り当てる。$W_i$ を $W$ のうち長さ $\le i$ のすべての元からなる加群とし、$\mathfrak a_i$ を $W_i$ における $\xi_i$ のすべての係数の系とする。まず次を得る。

$B$ が $W$ の真の約加群であり、$\mathfrak b_1,\ldots,\mathfrak b_k$ が $B$ に対応するイデアルであるなら、$\mathfrak b_i$ の中には、対応する $\mathfrak a_i$ の真の約イデアルが少なくとも一つある。$B\equiv0\pmod W$ から $B_i\equiv0\pmod {W_i}$ が従い、従って $i=1,2,\ldots,k$ に対して $\mathfrak b_i\equiv0\pmod{\mathfrak a_i}$ である。仮定により $B$ には $W$ に含まれない元が存在し、その中で最小の長さ $\ell$ をもつものがある。このとき $\mathfrak b_\ell$ は $\mathfrak a_\ell$ の真の約イデアルである。実際、
\[
        \beta=b_1\xi_1+\cdots+b_\ell\xi_\ell
\]
が $B$ におけるそのような最小長の元であるなら、$\mathfrak b_\ell\equiv0\pmod{\mathfrak a_\ell}$ である。もし $\mathfrak b_\ell=\mathfrak a_\ell$ なら、元
\[
        \alpha=a_\ell\xi_\ell+\cdots+a_1\xi_1\equiv0\pmod W
\]
が存在し、従って $\beta-\alpha\in B$, $\beta-\alpha\notin W$ で、長さが $\ell$ より小さい元が存在することになる。

いま約鎖または倍鎖
\[
        W^{(1)},W^{(2)},\ldots,W^{(\nu)},\ldots
\]
が与えられ、対応する鎖条件が $R$ で成り立つと仮定する。$W^{(\nu)}$ に割り当てられたイデアルを $\mathfrak a_{1\nu},\ldots,\mathfrak a_{k\nu}$ とすれば、列
\[
        \mathfrak a_{i1},\mathfrak a_{i2},\ldots
\]
は $i=1,\ldots,k$ について約イデアル鎖または倍イデアル鎖である。したがって仮定により、すべての $i$ について $\mathfrak a_{i\nu_0}$ が以後のすべてのイデアルと等しくなる添字 $\nu_0$ が存在する。すると、いま証明したことにより、加群 $W^{(\nu_0)}$ は以後のすべての加群と等しい。こうして鎖条件は移される。

これから直ちに次が得られる。

\medskip
\noindent\textbf{加群定理の帰結。} $R$ において、零イデアルと異なる任意のイデアルを法として倍イデアル鎖条件が成り立つなら、$M$ では、割り当てられたイデアル $\mathfrak c_1,\ldots,\mathfrak c_k$ がすべて零イデアルと異なる任意の加群 $C$ を法として倍加群鎖条件が成り立つ。これらの仮定の下で、倍イデアル鎖 $\mathfrak a_{i1},\ldots,\mathfrak a_{i\nu},\ldots$ は有限段階の後に停止する。

\medskip
\noindent\textbf{補足。} $R$ が単位元をもたない環である場合にも、零イデアルと異なるイデアルを法とする約イデアル鎖条件および倍イデアル鎖条件は移される。すなわち、$M$ の代わりに、形
\[
        a_1\xi_1+\cdots+a_k\xi_k+n_1\xi_1+\cdots+n_k\xi_k
\]
のすべての元からなる系 $M'$ を考える。この系は、追加の整倍 $n_i\xi_i$ を $M$ の元で置き換えると再び $M$ へ移る。この $M'$ に、上と対応して $2k$ 個のイデアルを割り当てることができる。ここで $\mathfrak a_i$ は $R$ のイデアルであり、$\mathfrak n_i$ は有理整数のイデアルである。$\mathfrak n_i$ については、約イデアル鎖条件と任意の非零イデアルを法とする倍イデアル鎖条件の双方が満たされるので、約イデアル鎖条件の証明は $M'$ に対して、従って $M$ に対しても有効である。しかし倍イデアル鎖条件については、$C$ あるいは $C'$ に割り当てられる $2k$ 個のイデアルがすべて零イデアルと異なることを要求しなければならない。

\subsection*{§3. 有限拡大体への移行}

数体と関数体の分類のためには、序で定式化した公理 I から V が、有限拡大体へ移る際にどのように移されるかを示す必要が残る。

1. \emph{$\mR$ において、公理 II、すなわち倍イデアル鎖条件を除くすべての公理 I から V が満たされているとする。$\mK$ を $\mR$ の商体とし、$\mL$ を $\mK$ 上の第一種有限拡大体とする。\footnote{Steinitz に従い、代数的拡大体は、そのすべての元が、適当な拡大体で相異なる線型因子に分解する素関数の零点であるとき、「第一種」と呼ばれる。} このとき、$\mR$ に関して整な $\mL$ のすべての元からなる系 $\mS$ では同じ公理が満たされ、$\mS$ の任意のオーダーでも整閉性を除くすべての公理がなお満たされる。}

§1, 3 により、$\mS$ は公理 III, IV、すなわち単位元の存在と零因子の不存在を満たす環をなす。§1, 4 の結論により、$\mS$ は $\mL$ において整閉である。同時に、$\mL$ は $\mS$ の商体になる。$\mL$ の任意の元は、$\mR$ の適当な元を掛けることで $\mS$ の元に移されるからである。したがって公理 V は満たされる。

I の証明はよく知られた議論に従う。第一種拡大として、$\mL$ は $\mK$ に一つの元 $\alpha$ を添加して得られ、この $\alpha$ は上により $\mR$ に関して整と仮定できる。$\alpha$ のほか、$\mK$ 上の $\mL$ の Galois 体に含まれる共役量 $\alpha',\alpha'',\ldots$ も整である。したがってそれらの差の積も整であり、その差の積の二乗 $D$ も整である。$\alpha,\alpha',\ldots$ はすべて異なるので、$D$ は $\mK$ の非零量であり、$\mK$ における $\mR$ の整閉性により $\mR$ の元である。通常の議論により、$\mR$ は整閉であるから、$\mS$ は元 $\xi_i=\alpha^i/D$ によって生成される有限 $\mR$-加群領域 $M$ に含まれる。このためには、$\mS$ の元を $\xi_i$ によって表した表示において、共役元へ移り、得られた線型方程式系を表示の係数について解けばよい。§2 の加群定理により、$M$ のすべての $\mR$-加群について約鎖条件が成り立ち、特に $\mS$ のすべてのイデアルについて成り立つ。$\mS$ の任意のオーダーのイデアルについても同様に約鎖条件が成り立つ。ここでも $M$ の中の $\mR$-加群を扱っているからである。さらに任意のオーダーについて、公理 III と IV は満たされる。

2. \emph{$\mR$ においてすべての公理 I から V が満たされるなら、$\mS$ においても同じことが成り立つ。$\mS$ の任意のオーダーでも、整閉性を除いてすべての公理がなお成り立つ。}

1 を踏まえれば、証明すべきは公理 II だけである。§2 の加群定理の系により、次を示せば十分である。$\mS$ の非零イデアルを $M$ の中の $\mR$-加群 $\mC$ と見たとき、$\mC$ に対応する $\mR$ のイデアル $\mc_i$ はすべて零イデアルと異なる。いま $\mC$ は、$\mK$ 上で $M$ と同じ有限線型階数をもつ。任意の元 $\gamma$ とともに、$\mC$ は $\gamma\alpha^i$ も含むからである。従って各 $\xi_i$ について、$c_i\xi_i$ が $\mC$ に属するような $\mR$ の元 $c_i\ne0$ が存在する。$\xi_i$ による表示の一意性により、ここで添字は体の次数より小さい値だけを走るものとするが、$c_i$ は $\mc_i$ に属し、従って $\mc_i$ は零イデアルではない。この議論は $M$ と同じ階数のすべてのオーダーに対して有効である。階数の小さいオーダーについてはその商体へ移る。この商体は $\mK$ と $\mL$ の間の中間体としてやはり第一種であり、$\mK$ 上の次数はオーダーの線型階数と一致する。従って同じ議論が保たれる。最後に、$\mS$ と異なる $\mS$ のオーダーは決して整閉ではないことを注意しておく。どのオーダーも $\mR$ を含むので、整閉性はそのオーダーに $\mS$ を含むことを強制するからである。

数体と関数体を従属させるには、したがって基礎領域、すなわち有理整数、一変数多項式、および複数不定元多項式の関数領域について、公理 I から V を検証すれば十分である。

3. \emph{有理整数の環 $\mR$、または体係数の一変数多項式の環については、公理 I から V が満たされる。} ここで I と II は、非零数または一変数の非零多項式を法とすると、剰余類が有限個、または線型独立な剰余類が有限個しかないという事実からも、またすべてのイデアルが主イデアルであるという事実からも従う。このことから、イデアルの素イデアル冪の積としての一意分解が従い、従って非零イデアルは有限個のイデアルによってしか割り切られない。公理 III と IV は明らかに満たされ、後者は等号の定義の帰結である。整閉性 V は、$\mR$ の元が単元、すなわち 1 の約元を除いて既約元へ一意に分解されることから従う。$\mR$ に含まれない商体の任意の元は、$\mR$ の二つの互いに素な元の商として書ける。従って分子のいかなる冪も分母で割り切られない。このような元は、$\mR$ に関して整であることを特徴づける方程式を満たし得ない。

4. \emph{いま $\mR$ を、体係数または有理整数係数の複数不定元の全多項式の環とする。} このとき公理 III, IV, V は 3 と同様に満たされる。ここでも元の既約元への因数分解は単元を除いて一意だからである。約イデアル鎖条件 I は、イデアル基底の存在に関する Hilbert の定理の直接の帰結であるが、倍イデアル鎖条件 II はここでは成り立たない。

しかし、\emph{関数領域}へ移る、すなわち最高次元のイデアルに制限することによって、公理 II の有効性も得ることができる。その場合、I の有効性も Hilbert の定理を用いずに 3 と同様に証明できる。実際、$\mR$ に不定元 $u$ を添加し、$u$ の全多項式で係数を $\mR$ にもつ環 $\mR^*$ を考える。等号は係数ごとに定義される。通常どおり、$\mR^*$ の多項式は、その係数の最大公約数、ただしイデアル約元でなく $\mR$ の多項式としての約元、が $\mR$ の単元であるとき、$\mR$ に関して原始的と呼ぶ。このとき $\mR^*$ のすべての多項式は、$\mR$ の元と $\mR^*$ の原始多項式の積であり、$\mR^*$ の原始多項式の積は再び原始である。

\emph{従って、$\mR^*$ の商体の元のうち、分母が $\mR^*$ の原始多項式であるもの全体は一つの環、すなわち $\mR$ の関数領域 $\mF$ をなす。} この関数領域 $\mF$ では、単元はちょうど、分子と分母がともに $\mR^*$ の原始多項式である元である。従って $\mF$ の各元は、$\mR$ の元と $\mF$ の単元の積として一意に表される。ゆえに $\mR$ の元に対する因数分解定理は $\mF$ の元へ移る。従って $\mF$ では、公理 III と IV に加えて、公理 V も 3 と同様に満たされる。

さらに $\mF$ ではすべてのイデアルが主イデアルである。なぜなら、イデアルは $\mF$ の任意の元とともに、それに単元を掛けて得られる $\mR$ の元も含み、また $\mR$ の任意の二つの元とともにその最大公約数も含むからである。実際、$f(x)=t(x)\overline f(x)$ および $g(x)=t(x)\overline g(x)$ とすれば、イデアルは $t(x)(\overline f(x)+u\overline g(x))$ も含み、$\overline f$ と $\overline g$ が互いに素でその線型結合が従って単元であるなら、$t(x)$ も含む。したがって、イデアルの任意の元 $f(x)$ から始めて、もしイデアルの中にそれで割り切れない元 $g(x)$ があれば、最大公約数 $t(x)$ もイデアルに属し、$f(x)$ の真の約元である。この操作を有限回繰り返すとイデアルの基底元に到達し、そのイデアルが主イデアルであることが分かる。従って $\mF$ では、$\mR$ の元の因数分解に対応して、イデアルの素イデアル冪の積としての一意分解が成り立つ。3 と同様に、公理 I と II が満たされることが従う。ここで考慮すべき $\mF$ の商体の拡大体は、$\mS$ の元を添加して生成できるものだけである。\footnote{これらの拡大体において $\mF$ に関して整な量の系は、$\mR$ から $\mF$ へ移ったのと同じように、$\mS$ から対応する関数領域へ移ることによっても得られる。J. König, \emph{Algebraische Größen}, Leipzig 1903, pp. 468/69 を参照。} 1 と 2 を考慮すれば、数体および関数体について公理 I から V の有効性が確立される。

\subsection*{§4. 同型定理。直和}

以下では、環性または加群性だけを仮定し、それ以外の公理は前提としない。

1. $M$ と $\overline M$ が $\mR$ に関する加群領域（§2）であるとき、$M$ が $\overline M$ に準同型（より正確には加群準同型）である、$M\sim\overline M$、とは、$M$ の各元に $\overline M$ の一つかつただ一つの元が対応し、$\overline M$ が尽くされ、\footnote{$\overline M$ で有効な等号定義の意味において。注6を参照。} かつこの対応の下で差と同じ $\mR$ の元による乗法が対応することをいう。従って $\beta\sim\overline\beta$ および $\gamma\sim\overline\gamma$ からは常に $(\beta-\gamma)\sim(\overline\beta-\overline\gamma)$ および $r\beta\sim r\overline\beta$ が従う。

従って $M$ の中の $\mR$-加群 $\mB$ には、$\overline M$ の中の $\mR$-加群 $\overline{\mB}$ が準同型的に対応する。また、$\overline M$ の $\mR$-加群 $\overline{\mC}$ の元に対応する $M$ の元全体 $\mC^*$ は、$\overline{\mC}$ によって一意に定まる $M$ の $\mR$-加群、すなわち $\overline{\mC}$ に随伴する加群 $\mC^*$ をなし、$\mC^*\sim\overline{\mC}$ である。特に $\mA$ が $\overline M$ の零元に随伴する $M$ の加群であるなら、$\mC^*$ は $\mA$ の約加群であり、$\mA$ を法として合同な $M$ の元の類に分かれる。それらの類は $\overline{\mC}$ の元と一対一に対応する。$M$ の中の $\mB$ から $\overline M$ の中の $\overline{\mB}$ へ移り、さらに戻って $\mB^*$ へ行くと、$\mB^*=(\mB,\mA)$、すなわち $\mB$ と $\mA$ の最大公約数を得る。従って $\mB$ が $\mA$ の約加群であれば $\mB^*=\mB$ である。

$M$ と $\overline M$ の元の対応が可逆的一対一であるなら、これらの領域は同型（より正確には加群同型）であると呼ばれ、$M\cong\overline M$ と書く。

$\mA$ が $M$ の任意の $\mR$-加群であるなら、$\mA$ を法とする合同を新しい等号関係とすることにより、$M$ に準同型な加群領域 $\overline M$、すなわち剰余類加群 $M\mid\mA$ が生じる。$M$ の各元には、$\overline M$ においてそれに等しいすべての、かつそれらだけの元が随伴する。$\overline M$ において等号定義から同一性へ移るとは、$\overline M$ で等しいすべての元を一つの類、すなわち剰余類にまとめ、これらの剰余類を $\overline M$ の新しい元として取ることである。

\emph{すべての準同型は剰余類加群への移行によって生じる。} 実際、$M\sim\overline M$ であり、$\mA$ が $\overline M$ の零元に随伴する加群なら、上で示したように $\overline M$ は剰余類加群 $M\mid\mA$ と同型である。

2. \textbf{第一同型定理。} \emph{$\overline M$ を剰余類加群 $M\mid\mA$ とし、$\mC$ を $\mA$ の約加群とする。このとき同型 $\overline M\mid\overline{\mC}\cong M\mid\mC$ がある。} 実際、$\mA$ を法とする合同は同時に $\mC$ を法とする合同であり、$\mA$ を法として等しい元は $\mC$ を法としても等しいままである。したがって、剰余類加群 $M\mid\mC$、すなわち $\mC$ を法とする等号は、まず $\mA$ を法として元を同一視し、$\overline M$ へ移り、ついで $\overline M$ の中で $\mC$ を法として等しい元をまとめることによって形成できる。これは $\overline M$ で $\overline{\mC}$ を法として同一視すること、すなわち $\overline M\mid\overline{\mC}$ を作ることと同じである。

\textbf{第二同型定理。} \emph{$\mB$ と $\mA$ が $M$ の加群であるなら、$(\mB,\mA)\mid\mA\cong\mB\mid[\mB,\mA]$ である。} なぜなら 1 により、$(\mB,\mA)\mid\mA$ を $\overline{\mB}$ と置くと、$\mB$ は $\overline{\mB}$ に準同型となる。そして $\overline{\mB}$ の零元に対応するのは正確に $[\mB,\mA]$ の元なので、主張した同型は再び 1 から従う。

3. $\mR$ と $\overline{\mR}$ が（可換）環であるとき、$\mR$ が $\overline{\mR}$ に準同型（より正確には環準同型）である、$\mR\sim\overline{\mR}$、とは、$\mR$ の各元に $\overline{\mR}$ の一つかつただ一つの元が対応して $\overline{\mR}$ が尽くされ、この対応の下で差と積が対応することをいう。元の対応が可逆的一対一であるなら、これらの環は同型（より正確には環同型）と呼ばれ、$\mR\cong\overline{\mR}$ と書く。

1 と 2 のすべての議論は、加群を $\mR$ のイデアルに置き換え、\footnote{非可換環では両側イデアルを用いなければならない。} 加群準同型および加群同型を環準同型および環同型に置き換えても有効である。特に、$\mc$ が $\ma$ の約イデアルなら、剰余類環 $\mc\mid\ma$ は、$\ma$ を法とする合同を新しい等号関係として導入することで得られる。ここでは、剰余類の積も定義されることに注意しなければならない。

従って $\mc$ は $\mc\mid\ma$ に準同型である。すべての準同型はこのようにして生じ、同型定理は 2 と同様に成り立つ。

\textbf{第一同型定理。} $\overline{\mR}$ が剰余類環 $\mR\mid\ma$ を表し、$\mc$ が $\ma$ の約イデアルであるなら、環同型の意味で $\overline{\mR}\mid\overline{\mc}\cong\mR\mid\mc$ である。

\textbf{第二同型定理。} $\mb$ と $\ma$ が $\mR$ のイデアルであるなら、環同型 $(\mb,\ma)\mid\ma\cong\mb\mid[\mb,\ma]$ が成り立つ。\footnote{群論でよく知られているこれらの同型定理は、加群についてはやや特殊な形で最初に Dedekind に現れる（第3版, p. 484 を参照）。上記のイデアルに対する形は Masazo Sono, \emph{On Congruences. I, II, III, IV}, Memoirs of the College of Science, Kyoto Imperial University, 2 (1917), 3 (1918, 1919), 2, p. 215 参照、に現れる。いずれの場合にも第二同型定理だけが明示的に述べられている。}

4. ここで、5 の直和に関する定理が従う互いに素なイデアルの計算法則\footnote{Dedekind（第4版, §178, III から VIII）に遡る形である。最近の叙述で一貫して現れる単位元を用いる計算は、はるかに煩雑である。}をまとめる。可換環 $\mR$ において、乗法単位元 $e$ の存在を仮定する。従って $\mR$ の任意のイデアルについて $\ma=\mo\ma$ であり、$\mo$ は単位イデアルを表す。二つのイデアル $\ma,\mb$ は、その最大公約数 $(\ma,\mb)$ が $\mo$ であるとき、互いに素と呼ばれる。

\emph{4$\alpha$. $(\ma,\mb)=\mo$ なら $(\ma,\mb\mc)=(\ma,\mc)$ である。} 実際、$(\ma,\mc)=(\ma,\mb)(\ma,\mc)=(\ma^2,\ma\mb,\ma\mc,\mb\mc)$ は $(\ma,\mb\mc)$ で割り切られ、逆も成り立つ。従って次が成り立つ。

\emph{4$\beta$. $\mc\mb\equiv0\pmod{\ma}$ かつ $(\ma,\mb)=\mo$ なら、$\mc\equiv0\pmod{\ma}$ である。} 実際、$\mc\mb\equiv0\pmod{\ma}$ は $(\ma,\mb\mc)=\ma$ と同値であり、$(\ma,\mb)=\mo$ であるから、$(\ma,\mc)=\ma$、すなわち $\mc\equiv0\pmod{\ma}$ も得られる。さらに次が成り立つ。

\emph{4$\gamma$. イデアル $\ma_1,\ldots,\ma_r$ の各々がイデアル $\mb_1,\ldots,\mb_s$ の各々と互いに素であるなら、$\ma_i$ の積は $\mb_i$ の積と互いに素である。} $(\ma,\mb)=\mo$ および $(\ma,\mc)=\mo$ から $(\ma,\mb\mc)=\mo$ が得られ、有限回繰り返せば主張が従う。

4$\alpha$ と 4$\gamma$ から次が従う。

\emph{4$\delta$. イデアル $\mb_1,\ldots,\mb_r$ が二つずつ互いに素、すなわち $i\ne k$ に対して $(\mb_i,\mb_k)=\mo$ であるなら、その最小公倍数は積に等しい。} $(\ma,\mb)=\mo$ とし、$\mv=[\ma,\mb]$ と置く。すると $\mv=\mo\mv=(\ma\mv,\mb\mv)=0(\ma,\mb)$ である。$\ma\mb\equiv0\pmod{\mv}$ だから、$\mv=\ma\mb$ が従う。4$\gamma$ を用いて有限回繰り返せば主張が得られる。

\emph{4$\varepsilon$. イデアル $\mb_1,\ldots,\mb_r$ が二つずつ互いに素であり、$\ma_i=[\mb_1,\ldots,\mb_{i-1},\mb_{i+1},\ldots,\mb_r]=\mb_1\cdots\mb_{i-1}\mb_{i+1}\cdots\mb_r$ と置くなら、$(\ma_1,\ldots,\ma_{i-1},\ma_{i+1},\ldots,\ma_r)=\mb_i$ であり、従って $(\ma_1,\ldots,\ma_r)=\mo$ である。} 実際、分配律により $(\ma_1,\ma_2)=\mb_3\cdots\mb_r(\mb_2,\mb_1)=\mb_3\cdots\mb_r$ である。従って $(\ma_1,\ma_2,\ma_3)=\mb_4\cdots\mb_r(\mb_3,\mb_1\mb_2)=\mb_4\cdots\mb_r$ であり、適切な番号づけの下で有限回繰り返せば主張が証明される。

イデアル $\ma_i$ は表示 $\mm=[\mb_1,\ldots,\mb_r]$ における $\mb_i$ の補イデアルと呼ばれる。

5. 再び、可換環 $\mR$ が乗法単位元をもつとする。$\mR$ の各元 $c$ が、$a_i$ を $\ma_i$ の元として、ただ一通りに $c=a_1+a_2+\cdots+a_r$ の形で表されるとき、$\mR$ はイデアル $\ma_1,\ldots,\ma_r$ の直和であるといい、記号で $\mR=\ma_1+\ma_2+\cdots+\ma_r$ と書く。

\emph{零イデアルが二つずつ互いに素なイデアル $\mb_1,\ldots,\mb_r$ の最小公倍数であり、$\ma_i$ が $\mb_i$ の補イデアルであるなら、$\mR$ はイデアル $\ma_1,\ldots,\ma_r$ の直和である。} $\mo=(\ma_1,\ldots,\ma_r)$ であるから、$\mR$ の任意の元は少なくとも一つの $\ma_i$ による加法表示をもつ。$[\ma_i,\mb_i]=(0)$ であるから、この表示は一意である。$0=a_1+a_2+\cdots+a_r=a_i+b_i$ から $a_i=0$ が得られる。第二同型定理により、$\ma_i$ は剰余類環 $\mR\mid\mb_i$ と同型である。直交関係 $i\ne k$ に対して $\ma_i\ma_k=(0)$、および $\ma_i^2=\ma_i$ が成り立つ。後者は $\ma_i=\mo\ma_i$ だからである。

\emph{逆に、$\mR$ の直和表示 $\mR=\ma_1+\ma_2+\cdots+\ma_r$ があり、$\mb_i=(\ma_1,\ldots,\ma_{i-1},\ma_{i+1},\ldots,\ma_r)=\ma_1+\cdots+\ma_{i-1}+\ma_{i+1}+\cdots+\ma_r$ と置くなら、$\mb_i$ は二つずつ互いに素であり、その最小公倍数は零イデアルである。} 実際、$\mo=(\ma_i,\mb_i)$ から、$k\ne i$ について $\mo=(\mb_k,\mb_i)$ が従う。$\mb_k$ は $\ma_i$ の約イデアルだからである。さらに $c$ がすべての $\mb_i$ で割り切れるなら、
\[
 c=a_2^{(1)}+\cdots+a_r^{(1)}=a_1^{(2)}+a_3^{(2)}+\cdots+a_r^{(2)}=\cdots=a_1^{(r)}+\cdots+a_{r-1}^{(r)},
\]
である。ただし各場合に $a_i^{(\lambda)}\equiv0\pmod{\ma_i}$ である。表示の一意性により、各表示で一つの成分が零であるから、すべての $\lambda$ について $0=a_1^{(\lambda)},\ldots,0=a_r^{(\lambda)}$ が従い、従って $c=0$ である。\footnote{5 に述べた定理は、直和と直交共通部分との関係に関する一般定理の特殊な場合である。H. Prüfer, \emph{Theorie der Abelschen Gruppen I}, Math. Zeitschr. 20 (1924), pp. 165--187, §6 を参照。そこに与えられた議論は、加群とイデアルが特殊な場合として含まれるほど一般的に群概念を定式化する場合にも有効である。}

\subsection*{§5. 素イデアルと準素イデアル}

再び $\mR$ を可換環とするが、ここではさらなる公理は仮定せず、単位元の存在さえ仮定しない。素イデアルと準素イデアルについての基本事実をまとめる。\footnote{\emph{イデアル論}, §4 では公理 I、すなわち約イデアル鎖条件を仮定している。そのため、素イデアルおよび準素イデアルのどの性質がこの公理に依存しないかは、そこでは鋭く切り分けられていない。}

1. $\mR$ のイデアル $\mpideal$ は、剰余類環 $\mR\mid\mpideal$ が零因子をもたない環であるとき、\emph{弱素イデアル}と呼ばれる。すなわち、$\ma\not\equiv0\pmod{\mpideal}$ および $\mb\not\equiv0\pmod{\mpideal}$ から常に $\ma\mb\not\equiv0\pmod{\mpideal}$ が従う場合である。$\mR$ のイデアル $\mpideal^*$ は、剰余類環 $\mR\mid\mpideal^*$ が零因子イデアルをもたない環であるとき、\emph{強素イデアル}と呼ばれる。すなわち、$\ma\not\equiv0\pmod{\mpideal^*}$ および $\mb\not\equiv0\pmod{\mpideal^*}$ から常に $\ma\mb\not\equiv0\pmod{\mpideal^*}$ が従う場合である。

\emph{すべての強素イデアルは同時に弱素イデアルである}。これは $\ma,\mb$ を主イデアルに特殊化すれば分かる。逆に、\emph{すべての弱素イデアルも強素イデアルである}。従って単に素イデアルといってよい。実際、$\mpideal$ を弱素イデアルとし、$\ma\not\equiv0\pmod{\mpideal}$ および $\mb\not\equiv0\pmod{\mpideal}$ とする。$\ma$ と $\mb$ の中から $a\not\equiv0\pmod{\mpideal}$ および $b\not\equiv0\pmod{\mpideal}$ を満たす元 $a,b$ を選ぶ。すると $ab\not\equiv0\pmod{\mpideal}$ であり、従って $\ma\mb\not\equiv0\pmod{\mpideal}$ であるから、$\mpideal$ は強素イデアルでもある。

2. $\mR$ のイデアル $\mq$ は、剰余類環 $\mR\mid\mq$ において、すべての零因子のある冪が消えるとき、\emph{弱準素イデアル}と呼ばれる。同値に、$a\not\equiv0\pmod{\mq}$ かつすべての $x$ について $b^x\not\equiv0\pmod{\mq}$ なら常に $ab\not\equiv0\pmod{\mq}$ である。$\mR\mid\mq$ で零因子となる $\mR$ のすべての元の系 $\mpideal$ はイデアルをなし、しかも素イデアルである。それは $\mq$ の約イデアルであり、随伴素イデアルと呼ばれる。実際、$a$ とともに $ra$ も零因子であり、$a$ と $b$ とともに $a-b$ も零因子である。$a^x$ と $b^\lambda$ があれば、$(a-b)^{x+\lambda}$ は常に $\mq$ で割り切られるからである。さらに $a$ が零因子でないなら、定義により $a$ のいかなる冪も零因子ではない。$a\not\equiv0\pmod{\mpideal}$ と $b\not\equiv0\pmod{\mpideal}$ から $a^x b^\lambda=(ab)^x\not\equiv0\pmod{\mq}$ が従い、従って $ab\not\equiv0\pmod{\mpideal}$ である。

$\mR$ のイデアル $\mq^*$ は、剰余類環 $\mR\mid\mq^*$ において、すべての零因子イデアルのある冪が消えるとき、\emph{強準素イデアル}と呼ばれる。すなわち、$\ma\not\equiv0\pmod{\mq^*}$ かつすべての $x$ について $\mb^x\not\equiv0\pmod{\mq^*}$ なら、常に $\ma\mb\not\equiv0\pmod{\mq^*}$ である。弱準素イデアルの場合と同様に、$\mR\mid\mq^*$ で零因子イデアルとなる $\mR$ のすべてのイデアルの最大公約数は、素イデアルであり、$\mq^*$ の約イデアルであり、随伴素イデアルをなすことが示される。逆に、同じ随伴素イデアル $\mpideal$ をもつすべての準素イデアルは、$\mpideal$ に属すると言われる。

\emph{すべての強準素イデアルは同時に弱準素イデアルである}。これは $\ma,\mb$ を主イデアルに特殊化すれば分かる。しかし一般には逆は偽である。\footnote{R. Hölzer が私に伝えた次の例がこれを示す。$\mR$ を体係数の可算個の不定元 $x_i$ の多項式領域とし、$\mq=(x_1^2,x_2^3,\ldots,x_\nu^{\nu+1},\ldots,x_i x_k\,[i\ne k]\ldots)$ とする。剰余類環における零因子はちょうど $\mpideal=(x_1,x_2,\ldots,x_\nu,\ldots)$ で割り切れる多項式であり、そのような各零因子のある冪は $\mq$ で割り切れる。従って $\mq$ は随伴素イデアル $\mpideal$ をもつ弱準素イデアルである。しかし $\mq$ は強準素イデアルではない。$\ma=(x_1,x_3,x_5,\ldots,x_{2\nu+1},\ldots)$ および $\mb=(x_2,x_4,\ldots,x_{2\nu},\ldots)$ と置くと、$\ma\mb\equiv0\pmod{\mq}$ であるが、$\ma$ または $\mb$ のどの冪も $\mq$ で割り切れない。} \emph{しかし $\mR$ に公理 I、すなわち約イデアル鎖条件を仮定すれば、すべての弱準素イデアルは強準素イデアルでもある。} 従って単に準素イデアルといってよい。$\mq$ を弱準素イデアルとし、$\ma\not\equiv0\pmod{\mq}$ かつすべての $x$ について $\mb^x\not\equiv0\pmod{\mq}$ とする。このとき、$a\not\equiv0\pmod{\mq}$ かつすべての $x$ について $b^x\not\equiv0\pmod{\mq}$ を満たす $\ma$ の元 $a$ と $\mb$ の元 $b$ が存在する。従って $ab\not\equiv0\pmod{\mq}$ であり、ゆえに $\ma\mb\not\equiv0\pmod{\mq}$ である。もし $\mb$ の各 $b$ について、ある指数 $\lambda$ が存在して $b^\lambda\equiv0\pmod{\mq}$ となるなら、有限個の基底元の指数を $\lambda_1,\ldots,\lambda_r$ として、$\mb^{\lambda_1+\cdots+\lambda_r}\equiv0\pmod{\mq}$ となるからである。\footnote{公理 I から有限イデアル基底を導くには、$\mR$ に固定された整列を用いなければならない。§6 を参照。} 同じ議論は、随伴素イデアルを $\mpideal$ とするとき、$\mpideal^\varrho\equiv0\pmod{\mq}$ となる最小指数 $\varrho$ が存在することも示す。この $\varrho$ は $\mq$ の指数と呼ばれる。

3. 以下では、約イデアル鎖条件を再び\emph{仮定しない}。有限個のイデアルの最小公倍数 $[\ma_1,\ldots,\ma_r]$ は、どの $\ma_i$ も残りのイデアルの最小公倍数に含まれない、すなわちどの $\ma_i$ も省けないとき、\emph{最短表示}と呼ばれる。

\emph{同じ素イデアル $\mpideal$ に属する有限個の弱準素イデアル、または強準素イデアルの最小公倍数は、再び $\mpideal$ を随伴素イデアルとする弱準素イデアルである。相異なる素イデアルに属する有限個の弱準素イデアル、または強準素イデアルによる最短表示は、弱準素イデアル、または強準素イデアルではない。} $\mf=[\mq_1,\ldots,\mq_r]$ とし、$\mpideal$ を弱準素イデアル $\mq_i$ の随伴素イデアルとする。このとき $\mpideal$ は、ある冪が $\mf$ で割り切れる元全体からも正確になる。$a\not\equiv0\pmod{\mf}$ かつすべての $x$ について $b^x\not\equiv0\pmod{\mf}$ なら、$b\not\equiv0\pmod{\mpideal}$ であり、また少なくとも一つの添字 $i$ について $a\not\equiv0\pmod{\mq_i}$ である。従って $ab\not\equiv0\pmod{\mq_i}$ であり、ゆえに $ab\not\equiv0\pmod{\mf}$ である。全体を通して元をイデアルに置き換えると、もはや $\mb\not\equiv0\pmod{\mpideal}$ を推論できない。

いま $\mm=[\mq_1,\ldots,\mq_r]$ を、$i\ne k$ なら $\mpideal_i\ne\mpideal_k$ である弱準素イデアルによる最短表示とし、$r\ge2$ とする。このとき少なくとも一つの随伴素イデアル、例えば $\mpideal_1$ は、他のどれにも含まれない。$\mpideal$ たちの任意の真の倍鎖は高々 $r$ 段で停止しなければならないからである。従って $\mpideal_i$ の元 $a_i$ で、$a_i^{\varrho_i}\equiv0\pmod{\mq_i}$ であり、かつ $a_2\not\equiv0\pmod{\mpideal_1},\ldots,a_r\not\equiv0\pmod{\mpideal_1}$ となるものが存在する。さらに $q_1\not\equiv0\pmod{\mm}$ を $\mq_1$ の元とすると、$q_1a_2^{\varrho_2}\cdots a_r^{\varrho_r}\equiv0\pmod{\mm}$ であるが、$a_2^{\varrho_2}\cdots a_r^{\varrho_r}$ のどの冪も $\mm$ で割り切れない。さもなければ $\mpideal_1$ による可除性が従うからである。従って $\mm$ は弱準素イデアルではない。全体を通して元をイデアルに置き換える、すなわち $q_1$ を $\mq_1$ に、$a_i$ を $\ma_i[\not\equiv0\pmod{\mpideal_1}$ だが $\ma_i^{\varrho_i}\equiv0\pmod{\mq_i}]$ に置き換えれば、強準素イデアルについての証明が得られる。\footnote{約イデアル鎖条件を不要にするこの私の元の証明の修正は、B. L. van der Waerden に負っている。}

4. \emph{加群商とイデアル商。} $\mT$ を $\mR$ の拡大環とし、$\mA,\mB,\mC,\ldots$ を $\mT$ の中の $\mR$-加群とする。商 $\mC=\mA:\mB$ は、条件 $\mC\mB\equiv0\pmod{\mA}$ と、もし $\mD\mB\equiv0\pmod{\mA}$ なら $\mD\equiv0\pmod{\mC}$ である、という条件を満たす加群として定義される。すなわち $\mC$ は、$\mC\mB\equiv0\pmod{\mA}$ が成り立つ最小の加群である。これは商を、$\mD\mB\equiv0\pmod{\mA}$ を満たすすべての加群 $\mD$ の最大公約数として一意に定める。そのような $\mD$ は存在する。零加群が確かにその一つだからである。

定義から次が従う。$\mB$ が $\overline{\mB}$ の倍加群であるなら、$\mA:\overline{\mB}$ は $\mA:\mB$ の倍加群である。$\mA$ が $\overline{\mA}$ の倍加群であるなら、$\mA:\mB$ は $\overline{\mA}:\mB$ の倍加群である。また
\[
        \mA:\mB\mC=(\mA:\mB):\mC=(\mA:\mC):\mB .
\]
拡大環 $\mT$ が $\mR$ と一致すれば、加群商はイデアル商 $\ma:\mb$ になる。従って上の規則もここで成り立つ。$\ma\mb\equiv0\pmod{\ma}$ であるから、$\ma\equiv0\pmod{\ma:\mb}$ でもある。

5. \emph{$\ma:\mb=\ma$ なら、$\mb$ は $\ma$ に対して素であると言う。} 従ってイデアル $\mb$ が $\ma$ に対して素であるための必要十分条件は、$\md\mb\equiv0\pmod{\ma}$ が常に $\md\equiv0\pmod{\ma}$ を含意することである。4 の計算法則から次が従う。$\mb$ と $\mc$ が $\ma$ に対して素なら、それらの積および最小公倍数も $\ma$ に対して素である。$\mb$ が $\ma$ に対して素で、$\ma$ が $\mb$ に対して素なら、$\ma$ と $\mb$ は互いに素と呼ばれる。§4, 4$\beta$ から、互いに素なイデアルは互いに素（相互素）でもあることが従う。


\setcounter{footnote}{25}

\subsection*{§6. 約イデアル鎖条件を仮定した場合のイデアル論}

基礎として、公理 I、すなわち約イデアル鎖条件を満たす可換環 $\mR$ を取る。以下ではさらに、$\mR$ の全元に一つの整列順序が固定されているものとする。これにより $\mR$ の全イデアルにも整列順序が定まる。実際、二つの仮定から、$\mR$ の各イデアルは有限個の元からなるイデアル基底をもつことが直ちに従う。$\mR$ の元の整列順序から、$\mR$ のすべての有限部分集合の準辞書式整列順序へ移り、
\footnote{これは次のように理解すべきである。有限部分集合 $\mathfrak W$ の元 $a_1,\ldots,a_n$ を、添字の順序が $\mR$ で指定された整列順序と一致するように書くとき、$n$ 個の元をもつ任意の部分集合は、$m>n$ 個の元をもつ部分集合に先行する。二つの部分集合が同じ個数の元をもつ場合は、整列されたリストが初めて異なる位置で、一方の該当元が他方の該当元に先行するとき、前者が後者に先行する。従って、有限部分集合の任意の系は第一元をもつ。$\mR$ の整列順序が与えられれば、それ以上の選択公理は不要である。とくに、代数的数体の場合のように $\mR$ が可算であれば、選択公理は関与しない。イデアル論における選択公理の役割は、詳述なしに \emph{イデアル論} 注 5 で述べておいた。}
各イデアルに対し、その可能な基底のうち、この有限部分集合の整列順序で第一のものを標準基底として割り当てるなら、有限部分集合の整列順序とともにイデアルの整列順序が得られる。

\textbf{定理 I.} \emph{約イデアル鎖条件の仮定の下で、$\mR$ の任意のイデアルは、有限個の既約イデアル、すなわち二つの真の約イデアルの最小公倍数として表せないイデアルの最小公倍数として表される。}

証明には次を示せば十分である。あるイデアル $\mm$ について定理 I が成り立たないなら、$\mm$ は、定理 I が同様に成り立たない真の約イデアルをもつ。このことから、約イデアル鎖条件に反して、有限段階で停止しない約イデアル鎖を構成できる。実際、$\mm$ は可約でなければならない。そうでなければ $\mm=[\mm]$ がすでに望む表示を与えるからである。$\mm=[\ma,\mb]$ とすると、定理 I が $\ma$ と $\mb$ の双方に同時に成り立つことはできない。もし成り立てば、$\mm$ についても対応する表示が従うからである。従って、定理 I が成り立たない $\mm$ の真の約イデアルが存在する。そのうち、イデアルの整列順序で第一のものを $\ma_1$ とする。同様に $\ma_1$ から真の約イデアル $\ma_2$ を構成し、一般に $\ma_i$ から真の約イデアル $\ma_{i+1}$ を構成すると、停止しない、しかもよく定まった約イデアル鎖が得られ、仮定に矛盾する。
\footnote{定理 I は、全加群の系に対して約加群鎖条件を仮定すれば、加群領域についてもまったく同様に成り立つ。この定理は純粋に集合論的な性格をもつ。同時に、証明においてイデアルの整列順序を用いる唯一の箇所でもある。抽象的には次の通りである。集合 $M$ において、部分集合の族 $\Sigma$ が区別され、整列されているとする。$\Sigma$-集合について鎖条件を仮定する。すなわち、$X_i$ が $X_{i-1}$ の真の上集合であるような鎖 $X_1,X_2,\ldots$ は必ず停止する。$\Sigma$-集合が、いずれも真の上集合である二つの $\Sigma$-集合の交わりであるとき可約といい、そうでないとき既約という。このとき任意の $\Sigma$-集合は、有限個の既約 $\Sigma$-集合の交わりである。}

約イデアル鎖条件を仮定しているので、§5, 2 により、単に準素イデアルと言うことができる。準素性と既約性の関係は次で与えられる。

\textbf{定理 II.} \emph{約イデアル鎖条件の仮定の下で、任意の既約イデアルは準素である。言い換えれば、準素でない任意のイデアルは可約である。}

剰余類環 $\mR/\mm$ へ移っても、約イデアル鎖条件は準同型により保たれる。$\mm$ の最小公倍数としての表示は、零イデアルの表示に対応し、また第一同型定理（§4, 3）により、$\mm$ の準素な約イデアルは $\mR/\mm$ の準素イデアルに対応し、逆も成り立つ。従って、剰余類環における零イデアルの分解を考えれば十分である。この環も同じ一般性をもつ環であるから、初めから分解すべきイデアルを $\mR$ の零イデアルと仮定してよい。

そこで、$\mR$ の零イデアルが準素でないと仮定する。少なくとも一対のイデアル $\ma,\mb$ が存在し、しかも $\mb$ は主イデアルと仮定してよく、次を満たす。
\[
\ma\ne(0),\qquad \mb^x\ne(0)\quad\text{任意の }x\text{ について},
\qquad\text{しかし}\qquad \ma\mb=(0).
\]
イデアル商の鎖
\[
\ma,\quad \ma:\mb,\quad \ldots,\quad \ma:\mb^\nu,\quad\ldots
\]
を作る。この鎖は停止する。例えば
\[
\mt=\ma:\mb^{m-1}=\ma:\mb^m=\cdots
\]
とする。従って $\mt=\mt:\mb$、すなわち $\mb$ は $\mt$ に対して素である。さらに $\mt$ は $\ma$ の約イデアルであって零イデアルと異なり、$\mb^x$ は任意の指数 $x$ について零と異なる。従って定理 II の証明には
\[
        (0)=[\mt,\mb^{m+1}]
\]
を示せば十分である。定義により
\[
        \mb^{m-1}\mt\equiv0\pmod{\ma},\qquad\text{従って}\qquad \mb^m\mt=(0),
\]
であるから、残るのは
\[
        \mc=[\mt,\mb^{m+1}]\equiv0\pmod{\mb^m\mt}
\]
を示すことだけである。これは、$\mb$ が主イデアルであり、$\mt$ に対して素であるという仮定から従う。$\mc$ の任意の元 $c$ は、$\mb^{m+1}$ で割り切れるので、$n$ を整数記号として
\[
        c=k b^{m+1}+n b^{m+1}=r b^m
\]
という表示をもつ。ただし $r$ は再び $\mR$ の元である。$c\equiv0\pmod{\mt}$ から $r b^m\equiv0\pmod{\mt}$ が従い、$\mt:\mb=\mt$ であるから $r\equiv0\pmod{\mt}$ が従う。よって $c\equiv0\pmod{\mt\mb^m}$ であり、定理 II は証明された。

\textbf{定理 III.} \emph{約イデアル鎖条件の仮定の下で、任意のイデアルは、相異なる素イデアルに属する有限個の最大準素成分の最小公倍数としての最短表示をもつ。}

このような表示を得るには、定理 I によって存在する既約イデアルによる表示、従って定理 II により準素イデアルによる表示を、可能な限り最短表示に置き換えればよい。同じ素イデアルに属する準素イデアルをまとめると、§5, 3 により主張された表示が得られる。これらの準素成分は最大と呼んでよい。なぜなら、それらのうち任意のものの最小公倍数はもはや準素ではないからである。
\footnote{対応する素イデアルと孤立成分は一意に定まる。\emph{イデアル論}、または W. Krull, ``Ein neuer Beweis für die Hauptsätze der allgemeinen Idealtheorie'', Math. Ann. 90 (1923), pp. 55--64 を参照。§7 では、そこに現れる単純な特殊場合について、一意性の直接証明を与える。そこに一意なものとして認識される準素イデアルは孤立成分である。}

\subsection*{§7. 二重鎖条件を仮定した場合のイデアル論}

これまで展開された理論に比べて得られる簡単化は、1 と 2 に掲げる補助結果に基づく。

1. \emph{零因子をもたない可換環で倍イデアル鎖条件が成り立つなら、その環はすでに体である。} 示すべきことは、$a\ne0$ に対して方程式 $ax=b$ が常に環の中に解をもつことである。解が高々一つであることは零因子がないことから従う。

$\ma$ を $a$ で生成される主イデアルとする。仮定により冪の列は停止する。$\ma^m$ が以後すべての冪に等しいとする。$\mb$ が $b$ で生成される主イデアルを表すなら、
\[
        \ma^m\mb=\ma^{m+1}\mb,
        \qquad \ma^{m+1}\mb=(a^{m+1}b).
\]
特に元 $a^m b$ は、ここでも $n$ を整数記号として、
\[
        a^m b=k a^{m+1}b+n a^{m+1}b=a^{m+1}c
\]
という表示をもつ。ただし $c$ は $\mR$ の元である。零因子がないので、これから $b=ac$ が従い、環が体であることが証明される。

2. 1 から直ちに次の補助結果が従う。\emph{倍イデアル鎖条件が可換環で成り立つなら、素イデアルは $\mo$ と異なる真の約イデアルをもたない。}

同様に、次の補足も得られる。\emph{公理 II だけ、すなわち零イデアルと異なる任意のイデアルを法とする倍イデアル鎖条件だけが満たされているなら、零イデアルと異なる任意の素イデアルは $\mo$ と異なる真の約イデアルをもたない。}

素イデアル $\mpideal$ を法とする剰余類環は、定義により零因子をもたない環である。倍イデアル鎖条件もそこで準同型により保たれるから、1 によりそれは体である。$\mpideal$ の約イデアルと剰余類環のイデアルとの一対一対応から、$\mpideal$ は $\mo$ 以外の真の約イデアルをもたない。
\footnote{$\mR$ に単位元が存在する必要はなく、従って $\mR$ の全元からなるイデアル $\mo$ に単位元が存在する必要もない。ただし 1 により、剰余類環には単位類が存在する。この種の最も単純な例は、補足の弱い場合では、すべての偶整数からなる系によって与えられる。}

\textbf{定理 IV.} \emph{可換環で二重鎖条件が成り立つなら、任意のイデアルの最短表示において、単位イデアル $\mo$ に属さない準素成分は一意に定まる。}

\emph{補足。} 公理 I と II の仮定の下で、定理 IV は零イデアルと異なる任意のイデアルについて成り立つ。

\[
        \mm=[\mq,\mq_1,\ldots,\mq_r]=[\mbar\mq,\mbar\mq_1,\ldots,\mbar\mq_s]
\]
を、最大準素成分による $\mm$ の最短表示とし、その随伴素イデアルを
\[
        \mo,\mpideal_1,\ldots,\mpideal_r,
        \qquad \mo,\mbar\mpideal_1,\ldots,\mbar\mpideal_s
\]
とする。$\mo$ に属する準素イデアルが現れないときは、成分 $\mq$、それぞれ $\mbar\mq$ を省くものとする。なぜなら
\[
        \mo\,\mpideal_1\cdots\mpideal_r\equiv0\pmod{\mm}
\]
であるから、各 $\mpideal_i$ はある $\mbar\mpideal_j$ に含まれなければならず、従って補助結果によりそれと同一である。逆に各 $\mbar\mpideal_j$ もある $\mpideal_i$ と一致しなければならない。従って $\mo$ と異なる素イデアルは一致する。さらに
\[
        \mq\,\mq_1\cdots\mq_r\equiv0\pmod{\mq_i}
\]
から $\mbar\mq_i\equiv0\pmod{\mq_i}$ が従う。なぜなら残りの成分は $\mpideal_i$ で割り切れず、従って $\mq_i$ に対して素だからである。逆向きの合同も同様に従うので、$\mo$ に属さない準素成分は一意である。$\mq$ が実際に現れるなら、表示が最短であるため $\mbar\mq$ も実際に現れなければならない。これにより随伴素イデアルの一意性も得られる。証明はまた、$\mo$ に属する準素成分を含まない任意の表示がすでに最短であることも示している。

3. 二重鎖条件にさらに単位元の存在を加えると、定理 IV は次の定理 V へと強められる。

\textbf{定理 V.} \emph{二重鎖条件を満たす単位元をもつ可換環では、任意のイデアルが、有限個の二つずつ互いに素な準素イデアルの積として一意に表される。}

\emph{補足。} 公理 I, II, III の下で、定理 V は零イデアルと異なる任意のイデアルについて成り立つ。

単位元が存在するので $\mo^2=\mo$ である。従って $\mo$ に属する任意の準素イデアルは $\mo$ に等しく、$\mo$ と異なるイデアルの最短表示には現れ得ない。従って定理 IV により準素成分は一意に定まり、同時に任意の表示は $\mo$ を省くことで最短表示となる。2 の補助結果により、相異なる二つの素イデアルは常に互いに素である。§4, 4 の計算規則により、その準素成分についても同じことが成り立つ。従って最小公倍数は積になる。

同じ仮定から次の補助結果が従う。\emph{単位元をもつ可換環で二重鎖条件が成り立つとき、イデアル $\mm$ に対して素な素イデアルは、単位イデアルを除けば、$\mm$ に含まれていないものにちょうど限られる。素であることと互いに素であることは一致する。}

\emph{補足。} 公理 I, II, III の下では、イデアル $\mm$ は零イデアルと異なるものと仮定しなければならない。

$\mpideal$ が $\mm$ に含まれていなければ、それは $\mm$ に属するすべての素イデアルと異なり、2 の補助結果によりそれらのいずれでも割り切れない。従って各準素成分に対して素であり、従って $\mm$ に対して素である。同時に各成分と互いに素であり、従って $\mm$ とも互いに素である。

他方、$\mpideal\ne\mo$ が $\mm$ に含まれていれば、それは随伴素イデアルの一つと一致しなければならない。それが成分 $\mq=\mq_1$ に属するとする。このとき $\mq:\mpideal$ は $\mq$ の真の約イデアルである。なぜなら
\[
        \mpideal^{\varrho-1}\equiv0\pmod{\mq:\mpideal},\qquad
        \mpideal^{\varrho-1}\not\equiv0\pmod{\mq}
\]
だからである。同時に $\mq:\mpideal$ は、$\mpideal^{\varrho-1}$ の約イデアルとして、素イデアル $\mpideal\ne\mo$ だけで割り切られ、従って準素である。分解の一意性により、積
\[
        (\mq:\mpideal)\mq_2\cdots\mq_r
\]
は $\mm:\mpideal$ で割り切れるが、$\mm$ の真の約イデアルである。従って $\mpideal$ は $\mm$ に対して素でない。

従って、イデアル $\mt$ が $\mm$ に対して素であるのは、$\mt$ に属するどの素イデアルも $\mm$ に含まれないとき、かつそのときに限る。この場合 $\mt$ は $\mm$ とも互いに素である。逆に、互いに素なイデアルは常に相互に素である（§5, 5）ので、この仮定の下では二つの概念は一致する。

\subsection*{§8. 商体における整閉性の下でのイデアル論}

ここまで展開されたイデアル論は、最後の二つの公理を加えることで通常の理論へと精密化される。

1. \emph{補助結果。} $\mR$ を零因子をもたず単位元をもつ可換環とする。$\mR$ に約イデアル鎖条件を仮定する（公理 I, III, IV）と、$\ma=\ma\mb$ かつ $\ma\ne(0)$ から $\mb=\mo$ が従う。従って
\[
        \mc\ne\mc\mt
\]
は、零イデアルおよび単位イデアルと異なるすべてのイデアルについて成り立つ。
\footnote{他方、同じ仮定の下でも、$\ma\mb=\ma\mc$ から $\mb=\mc$ を推論することはできない。例えば、$\mR$ を体上の $x,y$ の多項式領域とし、$\ma=(xy)$, $\mb=(x^3,xy,y^2)$, $\mc=(x^2,y^2)$ とする。このとき $\ma\mb=\ma\mc=\ma^2$ であるが、$\mb\ne\mc$ である。}
証明は補助結果 §1, 3 と同じである。$\ma\mb$ を $\mb$ に関する有限加群と見なすことができるからである。I により存在する $\ma$ のイデアル基底を $\alpha_1,\ldots,\alpha_n$ とする。$\ma=\ma\mb$ から方程式系
\[
        \alpha_i=b_{i1}\alpha_1+\cdots+b_{in}\alpha_n,
        \qquad b_{ij}\equiv0\pmod{\mb},\quad i=1,\ldots,n
\]
が得られる。$\mR$ は零因子をもたないので、行列式 $|b_{ij}-e_{ij}|$ は消える。ただし $i\ne j$ では $e_{ij}=0$、$e_{ii}=e$ とする。
\[
        e-b_{11}-\cdots\equiv0\pmod{\mb}
\]
から $e\equiv0\pmod{\mb}$、従って $\mb=\mo$ が従う。

2. あとは、公理 I から IV に商体における整閉性（公理 V）を加えると、零イデアルと異なる任意の準素イデアルがその随伴素イデアルの冪であることを示すだけである。結果として得られる表示
\[
        \mm=\mpideal_1^{\varrho_1}\cdots\mpideal_r^{\varrho_r}
\]
の一意性は、零イデアルと異なる任意のイデアルについて、すでに定理 V と 1 の補助結果により証明されている。同時に、これらの素イデアルは単位イデアルと異なる真の約イデアルをもたないことも示されている。証明は、まず Dedekind の帰結 II（§1, 4）を用いて指数二の場合に行い、一般の場合は完全帰納法によって行う。

\emph{補助結果。} \emph{公理 I から V の下で、指数二の準素イデアルは $\mq=\mpideal^2$ 以外に存在しない。}

示すべきことは、
\[
        \mpideal^2\equiv0\pmod{\mq},
        \qquad \mq\equiv0\pmod{\mpideal},
        \qquad \mq\not\equiv0\pmod{\mpideal^2}
\]
から必ず $\mq=\mpideal^2$ が従うことである。次を満たす $\mc$ を選ぶ。
\[
        \mc\equiv0\pmod{\mq},
        \qquad \mc\not\equiv0\pmod{\mpideal},
        \qquad \mc\equiv0\pmod{\mpideal^2}.
\]
補助結果 §7, 3 から、$\mc:\mpideal$ は $\mc$ の真の約イデアルになる。従って、ある元 $b$ が存在して
\[
        b\mpideal\equiv0\pmod{\mc}\equiv0\pmod{\mq},
        \qquad b\not\equiv0\pmod{\mc}
\]
を満たす。従って $y=b/c$ は商体に属するが $\mR$ には属さず、整ではない。示すべきは $b\not\equiv0\pmod{\mpideal}$ である。そうすれば、$\mq$ は準素で $\mpideal$ に属するので、$\mpideal\equiv0\pmod{\mq}$ が従う。

Dedekind の帰結 II により、$\mR$ の元 $m,n$ が存在して
\[
        y=b/c=m/n,
        \qquad m^2/n \text{ は整でない}
\]
となる。すなわち
\[
        bn=mc,
        \qquad m\not\equiv0\pmod{\mn},
        \qquad m^2\not\equiv0\pmod{\mn}
\]
である。$b\mpideal$ に $m$ を掛けると
\[
        mb\mpideal\equiv0\pmod{\mn c},
        \qquad\text{従って}\qquad m\mpideal\equiv0\pmod{\mn}
\]
を得る。従って $m\not\equiv0\pmod{\mpideal}$ である。なぜなら $m^2\not\equiv0\pmod{\mn}$ であり、また $\mn\not\equiv0\pmod{\mpideal}$ だからである。後者は、$\mn:\mpideal$ が $\mn$ の真の約イデアルであることから、再び §7, 3 により従う。四つの関係
\[
        m\not\equiv0\pmod{\mpideal},\qquad
        n\not\equiv0\pmod{\mpideal},\qquad
        c\equiv0\pmod{\mpideal},\qquad
        c\not\equiv0\pmod{\mpideal^2},\qquad
        bn=mc
\]
は今や $b\not\equiv0\pmod{\mpideal}$ を含意する。実際、$\mpideal^2$ は準素であるから、まず $mc\not\equiv0\pmod{\mpideal^2}$ が得られ、次に $bn=mc$ から $b\not\equiv0\pmod{\mpideal}$ が従う。これで補助結果が証明された。

3. \emph{補助結果。} \emph{公理 I から V の下で、指数 $\varrho$ の準素イデアルは $\mq=\mpideal^\varrho$ 以外に存在しない。}

これは 2 の補助結果から、さらなる公理の使用なしに従う。
\footnote{補助結果 2 の妥当性を仮定した補助結果 3 の証明は、Masazo Sono, \emph{On Congruences II}, §§9, 10 にある。注 17 を参照。}
まず次を示す。
\[
        \text{もし }\mc\equiv0\pmod{\mpideal},\ \mc\not\equiv0\pmod{\mpideal^2}
        \quad\text{ならば}\quad
        \mpideal^\varrho=(\mc^\varrho,\mpideal^{\varrho+1})
\]
がすべての $\varrho$ について成り立つ。2 の結果により $\mpideal=(\mc,\mpideal^2)$ である。$\mpideal^{\varrho-1}=(\mc^{\varrho-1},\mpideal^\varrho)$ と仮定して $\mpideal$ を掛け、分配法則を用いると、
\[
        \mpideal^\varrho=(\mc^\varrho,\mpideal^{\varrho+1})
\]
が得られる。求める結果は次の第二の形に置くことができる。もし
\[
        \mq\equiv0\pmod{\mpideal^\varrho},
        \qquad \mq\not\equiv0\pmod{\mpideal^{\varrho+1}},
        \qquad \mpideal^{\varrho+1}\equiv0\pmod{\mq},
        \quad \varrho\ge1,
\]
ならば、すでに $\mpideal^\varrho\equiv0\pmod{\mq}$ である。実際、任意の準素イデアルにはそのような指数 $\varrho>1$ がある。条件 $\mq\not\equiv0\pmod{\mpideal^{\varrho+1}}$ は、$\mpideal^\varrho\equiv0\pmod{\mq}$ と、1 の結果による $\mpideal^\varrho\ne\mpideal^{\varrho+1}$ から従う。第二の形を繰り返し適用すると $\mpideal^\varrho\equiv0\pmod{\mq}$ が得られ、従って $\mq=\mpideal^\varrho$ である。

第二の形の仮定と $\mpideal^\varrho$ の表示から、$\mq$ の各元 $q$ は
\[
        q=b c^\varrho+p^{\varrho+1},
        \qquad b\equiv0\pmod{\mpideal},
        \quad\text{または}\quad b c^\varrho\equiv0\pmod{\mq,\mpideal^{\varrho+1}}
\]
と書ける。$(\mq,\mpideal^{\varrho+1})$ は準素であるから、
\[
        c^\varrho\equiv0\pmod{\mq,\mpideal^{\varrho+1}},
        \qquad\text{または}\qquad
        c^{\varrho+1}\equiv0\pmod{\mq,\mpideal^{\varrho+2}}
\]
が従い、従って $\mpideal^\varrho\equiv0\pmod{\mq}$ である。

4. まとめると、次が得られる。

\textbf{定理 VI.} \emph{可換環で公理 I から V が成り立つなら、零イデアルおよび単位イデアルと異なる任意のイデアルは、零イデアルおよび単位イデアルと異なる有限個の素イデアルの冪の積として一意に表される。これらの素イデアルは、単位イデアルと異なる真の約イデアルをもたない。}

\subsection*{§9. 仮定された分解の帰結としての公理}

逆に、公理 I から V によって定理 VI から従う通常のイデアル分解の存在から、さらに公理を推論するためには、分解を次の形で仮定しなければならない。

\emph{仮定。} 可換環 $\mR$ において、零元または単位元によって生成されない任意のイデアルは、素イデアルの冪の積として一意に表される。これらの素イデアルは、単位元によって生成されない単純イデアル、すなわち $\mo$ と異なる真の約イデアルをもたないイデアルであり、逆にすべての単純イデアルは素イデアルである。一意性は、零元または単位元によって生成されない任意のイデアルについて
\[
        \ma^\varrho\ne\ma^{\varrho+1}
\]
が成り立つという鋭い形で成り立つものとする。

ここでは単位元の存在を仮定していない。単位元が存在しないなら、「単位元によって生成されない」という条件は制限ではない。

1. \emph{公理 III、単位元の存在の証明。} 公理 III が満たされていないと仮定する。仮定に反して、$\mo^2$ が素でない単純イデアルになることを示す。鋭い一意性仮定により $\mo\ne\mo^2$ である。従って $\mo\equiv0\pmod{\mo^2}$ だが $\mo^2\not\equiv0\pmod{\mo^2}$ であり、$\mo^2$ は素でない。しかしそれは単純である。$\mo$ と異なるどの素イデアルでも割り切れず、また仮定された積表示により、その約イデアルは $\mo$ と $\mo^2$ だけだからである。

2. \emph{公理 IV、零因子がないことの証明。} $a$ が零因子で、$a\ne0$, $b\ne0$ かつ $ab=0$ であるとする。すると $a$ と $b$ によって生成される主イデアルの積表示を掛け合わせて
\[
        (0)=\mpideal_1^{\varrho_1}\cdots\mpideal_r^{\varrho_r}
\]
を得る。ここで $\mpideal_i$ は零イデアルおよび単位イデアルと異なる。後者は、単位元の存在がすでに証明されているからである。表示は最短、すなわちどの因子を削除しても零イデアルの真の約イデアルを与えるものとして取る。零イデアルについては一意性を仮定していないので、これは自動的ではない。一つの素イデアルだけが現れるなら、$(0)=\mpideal^\varrho=\mpideal^{\varrho+1}$ となり、鋭い一意性要求に反する。複数の素イデアルが現れるなら、
\[
\begin{aligned}
\mpideal_1^{\varrho_1}
 &= (\mpideal_1^{\varrho_1},\mpideal_2^{\varrho_2}\cdots\mpideal_r^{\varrho_r})\mpideal_1^{\varrho_1}  \\
 &= \mpideal_1^{\varrho_1}\mpideal_2^{\varrho_2}\cdots\mpideal_r^{\varrho_r},
\end{aligned}
\]
である。なぜなら単位元の存在により、$\mpideal_1$ は他のすべての素イデアルと互いに素だからである。これも再び鋭い一意性条件に矛盾する。

3. \emph{公理 I と II、鎖条件の証明。} 仮定は、零イデアルと異なる任意のイデアルについて、可除性が積表示に反映されることを与える。従って、
\[
        \mm=\mpideal_1^{\varrho_1}\cdots\mpideal_r^{\varrho_r},
        \qquad
        \ma=\mbar\mpideal_1^{\bar\varrho_1}\cdots\mbar\mpideal_s^{\bar\varrho_s}
\]
で、$\mm\equiv0\pmod{\ma}$ かつ $\ma\ne\mo$ なら、各 $\mbar\mpideal_j$ はある $\mpideal_i$ と同一であり、$0\le\bar\varrho_i\le\varrho_i$ であることが従う。従って $\mm$ の約イデアルは、有限個の指数の組合せに対応して有限個しかない。よって、$\mm$ を法とする剰余類環において二重鎖条件が成り立つ。これにより公理 I と II が証明される。約イデアル鎖条件は $\mR$ 自身でも成り立つ。なぜなら零イデアルの真の約イデアルはすべて零イデアルと異なるからである。

公理 V、すなわち商体における整閉性の証明には、イデアル分解の標準的帰結を用いるので、まずそれらを導かなければならない。

4. \emph{零イデアルと異なる任意のイデアルを法とする剰余類環は主イデアル環である。} 剰余類環が零元だけからなる場合には主張は明らかなので、イデアルは単位イデアルと異なると仮定してよい。

まずイデアルが準素で、$\mq=\mpideal^\varrho$ とし、$\mc\equiv0\pmod{\mpideal}$ だが $\mc\not\equiv0\pmod{\mpideal^2}$ であるとする。このとき 3 により $\mc=\mpideal\ma$ で $(\ma,\mpideal)=\mo$ である。従って
\[
        \mpideal^\sigma=(\mc^\sigma,\mpideal^{\sigma+1})
        \qquad(\sigma<\varrho)
\]
であり、準素イデアルを法とする剰余類環の任意のイデアルは主イデアルである。一般に、
\[
        \mm=\mq_1\cdots\mq_r
\]
が表示であり、
\[
        \mR/\mm=\ma_1+\cdots+\ma_r
\]
が剰余類環の直和表示（§4, 5）であるなら、$\ma_i$ は $\mR/\mq_i$ と同型である。従って $\ma_i$ の任意のイデアルは主イデアルである。$\mc$ が剰余類環の任意のイデアルなら、$\ma_i\mc$ は主イデアル $\ma_i c_i$ であり、
\[
        \mc=\ma_1c_1+\cdots+\ma_r c_r=\mo c,
        \qquad c=c_1+\cdots+c_r.
\]
実際、$i\ne k$ で $\ma_i\ma_k=(0)$、かつ $c_i\equiv0\pmod{\ma_i}$ であるから、$\ma_i\mo c_i$ と $\ma_i\mc$ は $\ma_i$ の中で一致する。

剰余主イデアル環に関するこの定理は、次の形にも述べられる。$\mc$ が任意のイデアルなら、与えられたイデアル $\mb$ に対して素なイデアルを掛けることにより、$\mc$ を主イデアルに変えることができる。実際、$\mm=\mb\mc$ と置くと、$\mc$ は $\mm$ を法として主イデアルになる。すなわち
\[
        \mc=(\mo c,\mm)=(\mc\ma,\mc\mb)=\mc(\ma,\mb),
\]
従って $\mo c=\mc\ma$ かつ $(\ma,\mb)=\mo$ である。

5. \emph{分数イデアルの理論。} 分数イデアルとは、商体 $\mK$ 内の有限 $\mR$-加群を意味する。

\emph{$\mK$ 内の零でない有限 $\mR$-加群は、乗法に関してアーベル群をなす。}

二つの有限 $\mR$-加群の積は再び有限である。乗法は結合的かつ可換である。$\mK=\mo\mK$ であるから、単位イデアルが単位元である。残るのは、方程式
\[
        \mathfrak U X=\mo
\]
が有限 $\mR$-加群の系において常に解をもつことを示すだけである。

予備的注意。方程式 $\mathfrak A T=\mathfrak B$ がただ一つの解をもつなら、$T$ は加群商 $\mathfrak B:\mathfrak A$（§5, 4）である。なぜなら
\[
        T\equiv0\pmod{\mathfrak B:\mathfrak A},
        \qquad
        \mathfrak B=\mathfrak A T\equiv0\pmod{\mathfrak A(\mathfrak B:\mathfrak A)}\equiv0\pmod{\mathfrak B}
\]
だからである。

まず $\mathfrak A$ が主加群で、$\mathfrak A=(\alpha)$ であるなら、$X=(e/\alpha)$ は $\mathfrak A X=\mo$ の解であり、$X$ もまた有限 $\mR$-加群である。従って任意の有限 $\mR$-加群は、$\mR$ の二つのイデアルの加群商であり、このことが分数イデアルという名称を正当化する。実際、$t_1,\ldots,t_n$ が $T$ の加群基底であるなら、$t_i=\tau_i/a$ と置き、$\mt$ を $\mR$ において $\tau_i$ によって生成されるイデアルとする。このとき $\mo a\,T=\mt$ であり、予備的注意により $T=(\mt:\mo a)$ である。ここで商は $\mK$ の中で取られる。

方程式 $\mathfrak A X=\mo$ の解は一般に次のように与えられる。$\mathfrak A=\ma:\mo c$ と置き、$\ma\mb$ が主イデアル $\mo a$ になるように $\mb$ を選ぶ。このとき
\[
        \mathfrak A\,\mo c=\ma,
        \qquad \mathfrak A\,\mb\mo c=\mo a,
        \qquad \mathfrak A(\mb\mo c:\mo a)=\mo.
\]
ここで $X$ は、イデアルを主イデアルで割った商として、再び有限 $\mR$-加群である。群性が証明された。また任意の二つのイデアルの加群商 $\ma:\mb$ も、方程式 $\mb X=\ma$ の解として有限 $\mR$-加群であることが従う。

6. 群性からさらに次の帰結が従う。

6\,$\alpha$. \emph{消去が可能である。$\mathfrak T=\mathfrak A\mathfrak M:\mathfrak B\mathfrak M$ から $\mathfrak T=\mathfrak A:\mathfrak B$ が従い、逆も成り立つ。} 実際、$\mathfrak T\mathfrak B\mathfrak M=\mathfrak A\mathfrak M$ から $\mathfrak T\mathfrak B=\mathfrak A$ が従い、逆も同様である。

6\,$\beta$. \emph{任意の分数イデアルは、$\mathfrak C=\ma:\mb$ という表示をもつ。ただし $\ma$ と $\mb$ は互いに素であり、この条件により $\ma$ と $\mb$ は一意に定まる。} 存在は 5 と 6\,$\alpha$ から従う。$\mathfrak C=\ma:\mb=\mbar\ma:\mbar\mb$ から $\ma\mbar\mb=\mb\mbar\ma$ が従い、従って両方の組 $\ma,\mb$ および $\mbar\ma,\mbar\mb$ が互いに素であると仮定すれば一意性が得られる。

6\,$\gamma$. \emph{非整元、すなわち $\mR$ に属さない $\mK$ の元によって生成される任意の主加群は、分子のどの冪も分母で割り切れないような主イデアルの商として表示される。} 簡約表示を $\mo n=\mb:\mc$ とし、$(\mb,\mc)=\mo$ とする。4 に従って $\mo b=\mb\ma$ かつ $(\ma,\mc)=\mo$ となるように $\ma$ を選ぶ。すると
\[
        \mo n=\ma\mb:\ma\mc=\mo b:\ma\mc,
\]
であり、$\ma\mc=\mo b:\mo n$ なので、イデアル $\ma\mc$ もまた主イデアルである。従って
\[
        \mo n=\mo b:\mo c
\]
と書く。$(\ma,\mc)=\mo$ かつ $(\mb,\mc)=\mo$ であるから、$b$ のどの冪も $c$ で割り切れない。

7. \emph{公理 V、商体における整閉性の証明。} 6\,$\gamma$ により、$\mK$ の任意の非整元は商表示
\[
        \eta=a/c
\]
をもち、$\mR$ において分子のどの冪も分母で割り切れないようにできる。このような元は、$\mR$ 上整であることを特徴づける方程式を満たすことができない。なぜなら
\[
        \eta^n+r_1\eta^{n-1}+\cdots+r_n=0
\]
からは $\mR$ において $a^n\equiv0\pmod{\mo c}$ が得られるからである。

8. \emph{帰結。} 可換環 $\mR$ において公理 I から IV が成り立ち、かつすべての準素イデアルが既約であるなら、公理 V、すなわち商体における整閉性も成り立つ。公理 I から III により、任意の素イデアル冪 $\mpideal^\varrho$ は準素である。特に $\mpideal^2$ は準素であり、仮定により既約である。ここから
\[
        \mpideal=(\mo c,\mpideal^2)
\]
を証明しなければならない。そうすれば、補助結果 §8, 3 によりすべての準素イデアルは素イデアル冪となり、他の仮定と合わせて 7 から整閉性が従う。

約イデアル鎖条件により
\[
        \mpideal=(\mo c_1,\ldots,\mo c_n)
\]
と書ける。ここで乗数としては $\mpideal$ を法とする剰余類だけが関係し、$c_i$ は $\mpideal$ を法とする剰余類体上線型独立であると仮定できる。$n>1$ なら、この線型独立性から
\[
        \mpideal^2=[\mq_1,\mq_2],
        \qquad
        \mq_1=(\mo c_1,\mpideal^2),
        \qquad
        \mq_2=(\mo c_2,\ldots,\mo c_n,\mpideal^2),
\]
という表示が得られ、$\mq_1$ と $\mq_2$ はいずれも $\mpideal^2$ の真の約イデアルである。従って $\mpideal^2$ は可約となり、仮定に反する。よって $\mpideal=(\mo c,\mpideal^2)$ であり、主張した帰結が証明された。

逆に、公理 I から V の帰結として、任意の準素イデアルは直ちに既約である。素イデアル冪 $\mpideal^\varrho$ は、$\mpideal^\sigma$ および $\mpideal^\tau$ という形の真の約イデアルの最小公倍数にはなり得ないからである。

9. 環に零因子を許すなら、公理 I から III と商環における整閉性から、任意の準素イデアルが既約であることは従わない。
\footnote{注 18 を参照。}
$\mT$ を、整閉性を除く公理 I から IV がすべて成り立つ環とする。このような環は存在する。例えば、$\mR$ で公理 I から V が成り立つとき、$\mR$ の商体の有限拡大体における整量全体の系とは異なる有限オーダーがそれである（§3, 2）。8 の帰結により、$\mT$ の中には可約な準素イデアル $\mq$ が存在する。$\mpideal^\varrho$ を $\mq$ の真の倍イデアルとし、$\mR$ を剰余類環 $\mT/\mpideal^\varrho$ とする。この環では、$\mq$ に対応するイデアルが非零の可約準素イデアルである。$\mR$ では公理 I から III が成り立ち、商環における整閉性も成り立つ。なぜなら、$\mT$ の元で $\mpideal$ で割り切れないものはすべて $\mpideal^\varrho$ と互いに素になるので、$\mR$ のすべての正則元は単元であり、$\mR$ はその商環と一致するからである。

\subsection*{§10. 二重鎖条件と組成列}

ここで、整列されていると仮定された任意の加群領域（§2）について、二重鎖条件、すなわち加群の任意の約鎖および任意の倍鎖が有限段階で停止することが、組成列の存在と同値であることを示す。
\footnote{領域の加群は加法に関してアーベル群であり、正確には乗数領域が $\mR$ の元からなるので一般化されたアーベル群である。それらはアーベル群であるから、ここでは組成列の概念は主列の概念と一致する。この段落の定理は、「加群」の下に任意の非可換、さらには一般化された群の正規約数全体を理解しても成り立つ。以前の記法とわずかに異なる記法は、群論を想起させるためのものである。}
加群領域を可換環に特殊化すれば、環の全イデアルの系について対応する事実が得られる。2 では全体を通じて、加群同型を環同型に置き換えなければならない。

加群領域 $\mG$ は、$\mG$ 自身と零加群 $\mE$ 以外に $\mR$-加群を含まないとき、\emph{単純}と呼ばれる。加群 $\mA$ は、$\mG/\mA$ が単純であるとき、\emph{$\mG$ において単純}と呼ばれる。

倍鎖
\[
        \mG,\mU_1,\ldots,\mU_r,\mE
\]
は、鎖中のすべての加群が相異なり、各加群が直前のものにおいて単純であるとき、$\mG$ の長さ $r$ の組成列と呼ばれる。これは、剰余類加群、すなわち商群 $\mU_{i-1}/\mU_i$ が単純な加群領域であり、また $\mG/\mU_1$ と $\mU_r/\mE$ も単純であることと同値である。剰余類加群 $\mG/\mF$ を作ることにより、$\mG$ から $\mF\ne\mG$ まで組成列が存在する場合は、絶対的な場合に帰着される。

1. \emph{$\mG$ に二重鎖条件を仮定すると、$\mG\ne\mE$ なら $\mG$ に組成列が存在する。} $\mG$ の仮定された整列順序と約加群鎖条件から、§6 と同様に、準辞書式順序により全加群の系の整列順序が得られる。この整列順序と約加群鎖条件により、$\mG$ において単純な加群が少なくとも一つ存在する。すなわち、$\mG_1$ を整列順序で $\mG$ の第一の真の約加群とし、一般に $\mG_i$ を $\mG_{i-1}$ の第一の真の約加群とする。よく定まった鎖
\[
        \mG,\mG_1,\ldots,\mG_i,\ldots
\]
は有限段階で停止し、従って必ず $\mG$ において単純な加群 $\mA$ に到達する。$\mU_1\ne\mE$ なら、同じ手続きにより $\mU_1$ において単純な加群 $\mU_2$ が得られる。一般に、$\mU_{i-1}\ne\mE$ なら、$\mU_{i-1}$ において単純な加群 $\mU_i$ が存在する。これにより、よく定まった倍鎖
\[
        \mG,\mU_1,\ldots,\mU_i,
\]
が得られ、これは倍加群鎖条件により停止する。従ってこれは $\mG$ の組成列である。

2. \emph{$\mG$ に組成列が存在するなら、$\mG$ で二重鎖条件が成り立つ。} 証明は完全帰納法によるもので、$\mG$ の整列順序を必要としない。帰納法の過程で、組成列の存在だけを仮定した Jordan--Hölder の定理の証明も同時に得られる。
\footnote{環の場合については Masazo Sono（注 21）, \emph{On Congruences I}, §§11--13 を参照。そこでは非可換群における組成列の類似物も扱われている。すなわち $\mR,\mU_1,\ldots,\mU_r,\mE$ という列で、$\mU_i$ はイデアルであり $\mU_{i-1}$ において単純であるが、必ずしも $\mR$ のイデアルである必要はない。現在の議論は、倍鎖を、各 $\mU_i$ が $\mU_{i-1}$ において正規であるような鎖と解し、約鎖を対応して定義するなら、非可換群に対しても有効である。Dedekind, ``Über die von drei Moduln erzeugte Dualgruppe'', Math. Ann. 53 (1900), pp. 371--403 も参照。そこでは、はるかに一般的な領域について、存在仮定だけの下で同様に Jordan--Hölder の定理が証明されている。第二同型定理の代わりにやや弱い対応関係が現れるため、定理の主張もやや弱い。証明の原理はここに与えるものと同一である。}

$\mG$ が単純で、従って $r=0$ であり、唯一の組成列が $\mG,\mE$ であるとき、次の主張が成り立つ。

$\alpha$) $\mG$ において単純な任意の加群を通って組成列を引くことができる。

$\beta$) Jordan--Hölder の定理：任意の組成列は同じ長さをもち、商群（剰余類加群）の系は順序を除いて一致する。対応する商群は同型である。

$\gamma$) $\mG$ と異なる任意の加群を通って組成列を引くことができる。

$\delta$) 倍加群鎖条件が成り立つ。

$\varepsilon$) 約加群鎖条件が成り立つ。

従って、$\alpha$)--$\varepsilon$) が、長さ $r<r+1$ の組成列をもつ任意の加群領域について証明されていると仮定する。$\mG$ に長さ $r+1$ の組成列
\[
        \mG,\mA,\mA_1,\ldots,\mA_r,\mE
\]
が存在するという仮定の下で、それらを証明する。

2$\alpha$. $\mG$ において $\mA$ 以外に単純な加群がなければ、証明すべきことはない。$\mB$ を $\mG$ において単純な加群で、$\mB\ne\mA$ とする。$\mA$ と $\mB$ はいずれも $\mG$ において単純で、かつ異なるので、必ず $(\mA,\mB)=\mG$ である。$\mM=[\mA,\mB]$ と置く。第二同型定理（§4, 2）により
\[
        \mG/\mB\cong \mA/\mM,
        \qquad
        \mG/\mA\cong \mB/\mM.
\]
従って $\mM$ は $\mA$ においても $\mB$ においても単純である。$\mA$ は長さ $r<r+1$ の組成列をもつので、帰納仮定 $\alpha$ と $\delta$ により、$\mA$ の中に $\mM$ を通る組成列、例えば
\[
        \mA,\mM,\mM_1,\ldots,\mM_{r-1},\mE
\]
が存在する。従って
\[
        \mG,\mB,\mM,\mM_1,\ldots,\mM_{r-1},\mE
\]
は $\mB$ を通る $\mG$ の組成列である。

2$\beta$. Jordan--Hölder の定理を証明するため、
\[
\begin{aligned}
&(1)\quad \mG,\mA,\mA_1,\ldots,\mA_r,\mE,\qquad\text{および}\qquad
(2)\quad \mG,\mB,\mB_1,\ldots,\mB_s,\mE
\end{aligned}
\]
を $\mG$ の二つの組成列とする。$\mA=\mB$ なら、結果は長さ $r<r+1$ に対する帰納仮定から従う。そうでなければ、2$\alpha$ で構成された列
\[
        (3)\quad \mG,\mB,\mM,\mM_1,\ldots,\mM_{r-1},\mE,
        \qquad
        (4)\quad \mG,\mA,\mM,\mM_1,\ldots,\mM_{r-1},\mE
\]
と比較する。帰納仮定により、(1) と (3) の商群の系は一致する。同様に、適切な商へ移ると (4) は $\mB$ を通る長さ $r$ の組成列であるから、(2) と (4) の商群の系も一致する。従って $s=r$ である。2$\alpha$ の同型から (3) と (4) の一致が得られ、従って (1) と (2) の一致が得られる。これで Jordan--Hölder の定理が証明された。

2$\gamma$. 予備的注意。$\mA,\mB,\mH$ を $\mG$ の加群とし、$\mB$ が $\mA$ において単純であると仮定する。このとき加群 $(\mB,\mH)$ と $(\mA,\mH)$ は一致するか、あるいは $(\mB,\mH)$ が $(\mA,\mH)$ において単純である。第二同型定理により
\[
        (\mA,\mB,\mH)/(\mB,\mH)
        \cong
        \mA/[\mA,(\mB,\mH)]
\]
である。$\mB\equiv0\pmod{\mA}$ であるから、これは
\[
        (\mA,\mH)/(\mB,\mH)\cong \mA/\mC,
        \qquad \mC=[\mA,(\mB,\mH)]
\]
である。ここで $\mC\equiv0\pmod{\mA}$ である。他方、$\mB\equiv0\pmod{\mC}$ でもある。なぜなら $\mB\equiv0\pmod{\mA}$ かつ $\mB\equiv0\pmod{(\mB,\mH)}$ だからである。従って $\mC$ は $\mA$ と $\mB$ の間にあり、仮定により $\mA$ または $\mB$ と同一である。これで予備的注意は証明された。

今、
\[
        \mG,\mU,\mU_1,\ldots,\mU_r,\mE
\]
を $\mG$ の組成列とし、$\mH$ を $\mG$ と異なる任意の $\mG$ の加群とする。$\mH$ を通る組成列があることを示すため、
\[
        \mG,(\mU,\mH),(\mU_1,\mH),\ldots,(\mU_r,\mH),\mH
\]
を作る。予備的注意により、この中の相異なる加群は $\mG$ から $\mH$ への組成列をなす。従って第一の真の項は $\mG$ において単純である（これは、それが $\mB$ と一致する場合の 2$\alpha$ の特殊場合である）。2$\alpha$ により、$\mG$ にはその項を通る長さ $r<r+1$ の組成列があり、従って帰納仮定により $\mH$ を通る組成列がある。$\mG$ を付け加えれば、$\mH$ を通る $\mG$ の組成列が得られる。この議論を繰り返すと、そのような列はとくに、いま $\mG$ から $\mH$ まで構成した列の延長として選べることが分かる。

2$\delta$. 倍加群鎖条件の証明。再び $\mG$ に長さ $r+1$ の組成列が存在すると仮定し、
\[
        \mG,\mB,\mB_1,\ldots,\mB_i,\ldots
\]
を倍鎖とする。各項はその直前の項の真の倍加群である。鎖が $\mG$ から始まると仮定しても一般性を失わない。$\mG$ は常に第一項として付け加えられるからである。2$\gamma$ により、$\mB$ を通る組成列がある。従って $\mB$ はより短い長さの組成列をもつ。帰納仮定により、鎖
\[
        \mB,\mB_1,\ldots,\mB_i,
\]
は有限段階で停止する。従って $\mG$ を付け加えた後も同じことが成り立つ。

2$\varepsilon$. 約加群鎖条件の証明。再び $\mG$ に長さ $r+1$ の組成列を仮定し、
\[
        \mG,\mC,\mC_1,\ldots,\mC_i,
\]
を約鎖とする。各項はその直前の項の真の約加群である。この鎖が $\mG$ から始まると仮定しても、やはり制限はない。2$\gamma$ により、$\mC$ を通る組成列がある。従って $\mG/\mC$ にはより短い長さの組成列がある。帰納仮定により、約鎖
\[
        \mG/\mC,\ \mC_1/\mC,\ldots,\mC_i/\mC,
\]
は有限段階で停止する。第一同型定理（§4, 2）が与える一対一対応により、$\mG$ から始まる元の鎖についても同じことが成り立ち、従って $\mE$ を付け加えた鎖についても同じことが成り立つ。これにより、組成列の存在と二重鎖条件という二つの仮定の同値性が証明された。

\emph{1925年8月13日受理。}

\end{document}
